\documentclass[11pt,a4paper]{article}
\usepackage[T1]{fontenc}
\usepackage[utf8]{inputenc}
\usepackage{lmodern}
\usepackage{amsmath,amssymb,amsthm}
\usepackage[margin=1.15in]{geometry}
\usepackage[hidelinks]{hyperref}
\emergencystretch=3em

\newtheorem{theorem}{Theorem}[section]
\newtheorem{proposition}[theorem]{Proposition}
\newtheorem{lemma}[theorem]{Lemma}
\newtheorem{corollary}[theorem]{Corollary}
\theoremstyle{definition}
\newtheorem{definition}[theorem]{Definition}
\newtheorem{remark}[theorem]{Remark}
\newtheorem*{thmA}{Theorem A}
\newtheorem*{thmB}{Theorem B}
\newtheorem*{thmC}{Theorem C}
\newtheorem*{thmD}{Theorem D}

\newcommand{\C}{\mathbb{C}}
\newcommand{\Q}{\mathbb{Q}}
\newcommand{\Z}{\mathbb{Z}}
\newcommand{\F}{\mathbb{F}}
\newcommand{\Gm}{\mathbb{G}_m}
\newcommand{\one}{\mathbf{1}}
\newcommand{\Sym}{\operatorname{Sym}}
\newcommand{\Hilb}{\operatorname{Hilb}}
\newcommand{\tr}{\operatorname{tr}}
\newcommand{\ord}{\operatorname{ord}}
\newcommand{\Spec}{\operatorname{Spec}}
\newcommand{\Stab}{\operatorname{Stab}}
\newcommand{\Hom}{\operatorname{Hom}}
\newcommand{\diag}{\operatorname{diag}}
\newcommand{\rank}{\operatorname{rank}}
\newcommand{\Tr}{\operatorname{Tr}}

\title{Counterexamples to the positivity and saturation hypotheses of GCT~VI\thanks{AI-produced. The results are due to two runs of GPT-6
Astra (OpenAI, ChatGPT Pro mode); the text was written by Claude Opus~5.5
(Anthropic) at effort high, which also supplied the steps marked as added
in Section~\ref{sec:provenance}; see that section for who proved, wrote and
checked what. No human has checked this text. To the extent that any rights
exist in it, they are dedicated to the public domain under CC0~1.0.}}
\author{}
\date{25 November 2026}

\begin{document}
\maketitle

\begin{abstract}
Part VI of Mulmuley's geometric complexity theory formulates three hypotheses
on the stretching quasi-polynomials of Kronecker coefficients, plethysm
coefficients and, more generally, multiplicities of subgroup restriction:
the saturation index (SH) and the positivity index (PH2) should be bounded by
a polynomial in the size of the input, and the generating function should
admit a positive rational form whose largest denominator exponent, the modular
index, is bounded by such a polynomial (PH3). We show that these hypotheses
fail. For every $n\ge2$ the series $\sum_{s\ge0}g((s^n),(s^n),(s^n))\,t^s$
of triple-rectangular Kronecker coefficients has a pole at a primitive root of
unity of order $2^{n-1}3^{r}$ for some $r\ge0$, so every positive rational
form, indeed every expression $P(t)/\prod_i(1-t^{a_i})^{b_i}$, has some
$a_i\ge2^{n-1}$, while the input rank is $n^2+5n$; this refutes the existence
clause of PH3 for Kronecker coefficients, also in the relaxed form that
Mulmuley gave in the appendix to the paper of Briand, Orellana and Rosas. The
proof reduces modulo three to the cyclic permutation of the three tensor
factors, whose trace is the Hilbert series of the invariants of
$\C^n\oplus\mathbb S_{(2,1)}\C^n$, and locates there a unique fixed curve of
the grading group $\mu_{2^{n-1}}$. For the plethysm coefficients of
$\Sym^{d}(\Sym^3\C^n)$ along the ray of the determinant, the Hilbert series
has simple poles at the primitive $2^{n-1}$-th roots of unity for every
$n\ge4$, which refutes PH3; a signed trace of these poles dominates all other
terms, which gives positivity and saturation indices larger than $2^{cn}$ for
every $c<1/4$ and all large $n$, refuting PH2 and SH(a). Both arguments rest
on a support theorem for cubic forms: a diagonal matrix of two-power order
fixing a semistable cubic in $n$ variables has determinant of order at most
$2^{n-1}$, with equality only for the chain cubic
$x_1^2x_2+\dots+x_{n-1}^2x_n+x_n^3$. For restriction from a rank-one torus,
which GCT~VI admits as an explicit embedding, the three hypotheses fail at a
fixed input rank.
\end{abstract}

\tableofcontents

\section{Introduction and results}\label{sec:intro}

\paragraph{The hypotheses.}
Mulmuley's GCT~VI~\cite{GCT6} proposes to decide the nonvanishing of
structure constants of representation theory, among them the Kronecker
coefficients and the plethysm coefficients, by linear programming on a
stretching quasi-polynomial, and it formulates three hypotheses that would
make this work. Let $f(n)$ be the stretching function of a structure constant
(the constant at the $n$-fold stretched input), a quasi-polynomial, and
$F(t)=\sum_{n\ge0}f(n)t^n$ its generating function. The \emph{positivity
index} $p(f)$ and the \emph{saturation index} $s(f)$ measure how far $f$ must
be shifted before all constituents have nonnegative coefficients, respectively
are positive at every positive integer whenever not identically zero; a
\emph{positive form} of $F$ is an expression with denominator
$\prod_i(1-t^{a_i})^{d_i}$ and a numerator with nonnegative integer
coefficients and constant term one, and its \emph{modular index} is
$\max_ia_i$. (The definitions are recalled in Section~\ref{sec:defs}.) The
hypotheses ask that $s(f)$ (SH(a)), $p(f)$ (PH2) and the modular index of some
positive form (PH3) be bounded by a polynomial in the rank of the input:
\cite[Hypotheses 3.3.2(a), 3.3.3, 3.3.4, p.~51]{GCT6} in general and
\cite[Hypotheses 1.6.5(a), 1.6.6, 1.6.7, pp.~24--25]{GCT6} for plethysm.
After Briand, Orellana and Rosas found counterexamples to the earlier, strict
forms of SH and PH2 for Kronecker coefficients~\cite{BOR}, Mulmuley stated
relaxed forms of SH, PH2 and PH3 for Kronecker coefficients in an appendix to
their paper~\cite[Appendix, pp.~18--20]{BOR}; the polynomial bounds above are
of that relaxed kind.

\paragraph{Results.}
Write $g(\lambda,\mu,\nu)$ for the Kronecker coefficient and put
\[
 k_n(s)=g\bigl((s^n),(s^n),(s^n)\bigr),\qquad
 K_n(t)=\sum_{s\ge0}k_n(s)t^s .
\]
This is the stretching function of the Kronecker coefficient of the triple of
columns $((1^n),(1^n),(1^n))$.

\begin{thmA}[Theorem~\ref{thm:kronecker}]
For every $n\ge2$ there is an integer $r\ge0$ such that $K_n(t)$ has a pole at
every primitive root of unity of order $2^{n-1}3^r$. Consequently every
expression $K_n(t)=P(t)/\prod_i(1-t^{a_i})^{b_i}$ with $P\in\Q[t]$ and
positive integers $a_i,b_i$ has $\max_ia_i\ge2^{n-1}$. The existence clause of
PH3 for Kronecker coefficients \cite[Hypothesis 3.3.4]{GCT6} fails, and so
does the relaxed PH3 of \cite[Appendix]{BOR}.
\end{thmA}

Let $\nu=\gcd(n,3)$, $d_0=n/\nu$ and $m=\binom{n+2}3=\dim\Sym^3\C^n$. The
plethysm coefficient of the ray $\mu=(3)$, $\lambda=(d_0)$, $\pi=((3/\nu)^n)$
of \cite[Section 1.6]{GCT6} has stretching function
\[
 f_n(s)=\dim\Hom_{GL_n}\bigl((\det)^{3s/\nu},\Sym^{d_0s}(\Sym^3\C^n)\bigr),
 \qquad H_n(t)=\sum_{s\ge0}f_n(s)t^s ,
\]
the number of linearly independent $SL_n$-invariants of degree $d_0s$ of
cubic forms in $n$ variables.

\begin{thmB}[Theorem~\ref{thm:plethysmPH3}]
For every $n\ge4$, $H_n(t)$ has a simple pole at every primitive root of unity
of order $2^{n-1}$. Every expression $H_n(t)=P(t)/\prod_i(1-t^{a_i})^{b_i}$ has
some $a_i$ divisible by $2^{n-1}$, while the parameter of
\cite[Hypothesis 1.6.7]{GCT6} is $m^2+n+1=O(n^6)$. PH3 fails for plethysm
coefficients.
\end{thmB}

\begin{thmC}[Theorem~\ref{thm:indices}]
For every real $0<c<1/4$ and all sufficiently large $n$, the positivity index
and the saturation index of $f_n$ satisfy
\[
 p(f_n)>2^{cn},\qquad s(f_n)\ge2^{cn};
\]
explicitly, $p(f_n)>2^{n/8}$ and $s(f_n)\ge2^{n/8}$ for $n\ge2\cdot10^{80}$.
PH2 and SH(a) for plethysm coefficients \cite[Hypotheses 1.6.5(a),
1.6.6]{GCT6} fail.
\end{thmC}

\begin{thmD}[Propositions~\ref{prop:torusmodular} and~\ref{prop:torusindices}]
For restriction from $\Gm$ along explicit polynomial embeddings into $GL_2$
and $GL_3$, which \cite[Section 3.4]{GCT6} admits, the input rank stays $5$,
respectively $6$, while the modular index of every positive form, respectively
the positivity and saturation indices, are unbounded. Hypotheses 3.3.2(a),
3.3.3 and 3.3.4 of \cite{GCT6} fail as stated for the subgroup restriction
problem.
\end{thmD}

\paragraph{What is not claimed.}
None of these results bears on the density clauses \cite[Hypotheses
1.6.5(b), 3.3.2(b)]{GCT6}, which ask for strict saturation of almost all
inputs; our families are sparse. We do not determine the exponent $r$ of
Theorem~A (for $n=2,3,4,5$ separate computations in the run give $r=0$; we do
not use them), and we do not bound the positivity or saturation index of the
Kronecker family. Nothing here concerns the occurrence of representations in
orbit closures, the other route of geometric complexity theory, or the
complexity of deciding the vanishing of these coefficients.

\paragraph{Method.}
The refutations of PH3 for Kronecker and for plethysm coefficients detect a pole of a Hilbert series at a root of
unity $\zeta$ of large order by the fixed locus of the grading element
$\zeta$ on the invariant-theoretic quotient: a unique fixed curve along which
the quotient is smooth forces a pole (Proposition~\ref{prop:fixedcomponentpole}),
while an isolated fixed point forces regularity (Lemma~\ref{lem:isolated}).
The fixed loci are found by a support theorem for cubic forms
(Theorem~\ref{thm:fullcubicsupport}), which singles out the chain cubic. For
Kronecker coefficients the quotient is that of three-tensors, whose fixed loci
we cannot control directly; instead the cyclic permutation $\tau$ of the three
factors acts on the invariants, its trace series agrees with $K_n$ modulo
three, and by a character identity (Theorem~\ref{thm:K023}) it is the Hilbert
series of a quotient of $\C^n\oplus\mathbb S_{(2,1)}\C^n$, where the support
theorem applies again. A transfer lemma for congruent rational series
(Theorem~\ref{thm:transfer}) carries the pole back, at the price of the unknown power $3^r$.
The torus examples of Theorem~D are elementary.

\paragraph{Earlier work.}
On 25 November 2026 we ran web searches, searched the Hexagon repository, and
went through the papers that cite \cite{GCT6} or \cite{BOR} in the OpenAlex
index, for earlier statements that PH3, PH2 or SH fails in these forms;
Section~\ref{sec:literature} lists what was searched and what was not.
Briand, Orellana and Rosas~\cite{BOR} refuted the strict forms of SH and PH2
for Kronecker coefficients; the positivity indices of their examples are at
most one, so they agree with the relaxed forms. Ikenmeyer, Mulmuley and
Walter~\cite{IMW} and Fischer and Ikenmeyer~\cite{FI} showed that deciding the
positivity of Kronecker and of plethysm coefficients is NP-hard; this concerns
the exact decision problem, not the index hypotheses. Kahle and
Micha\l{}ek~\cite{KM} showed that plethysm quasi-polynomials need not count
lattice points of polytopes of the matching dimension, which bears on the
polytope hypothesis PH1. In these searches we found no earlier refutation of
PH3, or of the relaxed SH(a) and PH2.

Section~\ref{sec:provenance} says who proved and checked what.

\section{Definitions and the hypotheses}\label{sec:defs}

\paragraph{Quasi-polynomials and indices.}
A function $f:\Z_{\ge0}\to\Z$ is a quasi-polynomial if there are a period
$l\ge1$ and polynomials $f_1,\dots,f_l$ with $f(n)=f_j(n)$ for $n\equiv j
\pmod l$. The following are \cite[Definitions 1.2.2 and 1.2.4,
pp.~13--14]{GCT6}.

\begin{definition}\label{def:indices}
The quasi-polynomial $f$ is \emph{strictly positive} if all coefficients of
all $f_j$ are nonnegative; its \emph{positivity index} $p(f)$ is the least
integer $p\ge0$ such that $n\mapsto f(n+p)$ is strictly positive. It is
\emph{strictly saturated} if, for every $j$, either $f_j$ is identically zero
or $f_j(n)>0$ for every integer $n\ge1$; its \emph{saturation index} $s(f)$ is
the least integer $s\ge0$ such that $n\mapsto f(n+s)$ is strictly saturated.
\end{definition}

\begin{definition}[{\cite[Definition 1.2.7, p.~15]{GCT6}}]\label{def:positiveform}
Let $F(t)=\sum_{n\ge0}f(n)t^n$ with $f$ a quasi-polynomial of degree $d$, not
identically zero. A \emph{positive form} of $F$ is an expression
\[
 F(t)=\frac{h_0+h_1t+\dots+h_et^e}{\prod_{i=0}^k(1-t^{a_i})^{d_i}}
\]
with $h_0=1$, all $h_i$ nonnegative integers, $a_i,d_i$ positive integers and
$\sum_id_i=d+1$. Its \emph{modular index} is $\max_ia_i$.
\end{definition}

In \cite[eq.~(1.4)]{GCT6} the numerator is printed as $h_dt^d+\dots+h_0$,
with the letter $d$ of the degree of $f$. We allow a numerator of any degree
$e$; this only enlarges the class of positive forms.

Our lower bounds on modular indices use only the denominator: they hold for
every expression $P(t)/\prod_i(1-t^{a_i})^{b_i}$ with $P\in\C[t]$, with no
condition on the numerator or on the exponents $b_i$.

\paragraph{The hypotheses.}
In \cite[Section 3.3, p.~51]{GCT6} the input $x$ of a structure constant has a
rank $\rank(x)$ and a stretching quasi-polynomial $\tilde f(x,n)$, and the
hypotheses read: (SH)(a) $s(\tilde f)=O(\mathrm{poly}(\rank(x)))$, (b) the
saturation index vanishes for almost all $x$ (Hypothesis 3.3.2); (PH2)
$p(\tilde f)=O(\mathrm{poly}(\rank(x)))$ (Hypothesis 3.3.3); (PH3) the
generating function $F(x,t)=\sum_n\tilde f(x,n)t^n$ has a positive rational
form of modular index $O(\mathrm{poly}(\rank(x)))$, and more specifically the
modular index defined in \cite[Section 4.1.1]{GCT6} is
$O(\mathrm{poly}(\rank(x)))$ (Hypothesis 3.3.4). These are stated for the
Kronecker problem, the plethysm problem and the subgroup restriction problem,
among others. For subgroup restriction along an explicit polynomial
homomorphism $\rho:H\to G$, with $H$-representation $\lambda$ and
$G$-representation $\pi$, the rank is the sum of the ranks of $H$ and $G$ and
the heights of the specified weights \cite[Section 3.4, p.~52]{GCT6}; the
Kronecker and plethysm problems are special cases.

For plethysm, \cite[Section 1.6, pp.~24--25]{GCT6} states the same hypotheses
with the parameter $h=\dim G+\operatorname{ht}(\lambda)+\operatorname{ht}(\pi)$:
SH(a) asks $s\le ch^{c'}$ (Hypothesis 1.6.5(a)), PH2 asks for a polynomial
bound on the positivity index (Hypothesis 1.6.6), and PH3 asks that the
generating function have a positive rational form with modular index bounded
by a polynomial in $h$, and more specifically that this hold for the minimal
positive form attached to an h.s.o.p.\ of lexicographically minimal degree
sequence (Hypothesis 1.6.7).

The relaxed hypotheses of \cite[Appendix, pp.~19--20]{BOR} are stated for
Kronecker coefficients $g_{\lambda,\mu}^{\pi}$ with $H=GL_n\times GL_n\subset
G=GL_{n^2}$, with the parameter $h$ the largest height of $\lambda,\mu,\pi$:
(SH)(a) $s=O(\mathrm{poly}(h))$; (PH2) $p=O(\mathrm{poly}(h))$; (PH3) the
modular index of the minimal positive form is $O(\mathrm{poly}(h))$.

\paragraph{The inputs of this paper.}
For $k_n$ we take $H=GL_n\times GL_n$, $G=GL_{n^2}$ and the columns
$\lambda=\mu=(1^n)$, $\pi=(1^n)$, stretched to $(s^n)$. Then $\rank(x)=2n+n^2+3n
=n^2+5n$ in the sense of \cite[Section 3.4]{GCT6}, and $h=n$ in the sense of
\cite[Appendix]{BOR}. For $f_n$, $G=GL(\Sym^3\C^n)=GL_m$, $\operatorname{ht}
(\lambda)=1$, $\operatorname{ht}(\pi)=n$, so $h=m^2+n+1$, and the rank of
\cite[Section 3.4]{GCT6} is $n+m+n+1$. Both are bounded by a polynomial in
$n$; so is the parameter of every hypothesis above.

\section{Hilbert series and fixed points of the grading}\label{sec:hilbert}

Throughout, $R=\bigoplus_{d\ge0}R_d$ is a finitely generated positively graded
$\C$-algebra with $R_0=\C$, and $\Hilb_R(t)=\sum_d\dim R_d\,t^d$. It is a
rational function whose poles are roots of unity. For a root of unity $\eta$,
the grading element $\eta$ acts on $R_d$ by $\eta^d$; we write $X=\Spec R$ and
$X^\eta$ for the scheme-theoretic fixed locus. If $\eta$ has order $q$, then
$X^\eta=\Spec R/I_q$ with
\begin{equation}\label{eq:fixedideal}
 I_q=(R_d:\ q\nmid d),
\end{equation}
since a homogeneous element of degree $d$ is changed by the factor $\eta^d$.

\begin{lemma}\label{lem:isolated}
If $X^\eta$ is supported at the graded origin, then $\Hilb_R(t)$ is regular at
every primitive root of unity of the order $q$ of $\eta$.
\end{lemma}

\begin{proof}
Choose homogeneous generators $a_1,\dots,a_l$ of $I_q$ with degrees $e_i$ not
divisible by $q$. As the support is the origin, $R/I_q$ is finite-dimensional.
The homology modules $N_j$ of the graded Koszul complex $K_R(a_1,\dots,a_l)$
are finitely generated over $R$ and annihilated by $I_q$, because
multiplication by each $a_i$ is null-homotopic on the Koszul complex; hence
they are finite-dimensional. The graded Euler characteristic gives
\[
 \Hilb_R(t)\prod_{i=1}^l(1-t^{e_i})=\sum_j(-1)^j\Hilb_{N_j}(t),
\]
whose right side is a polynomial. No factor $1-t^{e_i}$ vanishes at a
primitive $q$-th root of unity.
\end{proof}

\begin{proposition}\label{prop:fixedcomponentpole}
Let $R$ be a domain, $M\ge2$, $\zeta$ a primitive $M$-th root of unity. Suppose
$X^\zeta$ has a unique irreducible component $Y$ of maximal dimension
$\delta\ge1$, and $X$ is smooth at the generic point of $Y$. Give $Y$ its
reduced structure; its coordinate ring is concentrated in degrees divisible
by $M$; put $h_Y(z)=\sum_k\dim\C[Y]_{Mk}z^k$ and
$e_Y=\lim_{z\to1}(1-z)^\delta h_Y(z)>0$. Then the conormal space of $Y$ in $X$
at its generic point has no eigenvalue one under $\zeta$, and if
$\nu_1,\dots,\nu_c$ ($c=\dim X-\delta$) are its eigenvalues,
\begin{equation}\label{eq:residue}
 \lim_{t\to\zeta}(1-t^M)^\delta\Hilb_R(t)=\frac{e_Y}{\prod_{j=1}^c(1-\nu_j)}\ne0 .
\end{equation}
In particular $\Hilb_R$ has a pole of order exactly $\delta$ at $\zeta$, and
every expression $\Hilb_R(t)=P(t)/\prod_i(1-t^{a_i})^{b_i}$ with $P\in\C[t]$
has some $a_i$ divisible by $M$.
\end{proposition}

\begin{proof}
Choose homogeneous algebra generators of $R$ and list those of degree not
divisible by $M$ as $f_1,\dots,f_l$, $\deg f_i=w_i$. Their ideal $I$ is the
fixed-locus ideal: $f_i$ is a nonzero multiple of $\zeta\cdot f_i-f_i$, and
modulo the $f_i$ every remaining generator is fixed. Thus $X^\zeta=\Spec R/I$
and $(R/I)_d=0$ when $M\nmid d$. The ideal is homogeneous, so every component
of the fixed locus is stable under the connected grading torus, hence graded;
the coordinate ring of $Y$ lives in degrees divisible by $M$.

Let $K_\bullet$ be the graded Koszul complex on $f_1,\dots,f_l$, its $i$-th
exterior generator of degree $w_i$ and $\zeta$-character $\zeta^{w_i}$, and
$N_j$ its homology. As before, each $N_j$ is a finite graded $R/I$-module, and
\begin{equation}\label{eq:euler}
 \Hilb_R(t)\prod_{i=1}^l(1-t^{w_i})=\sum_{j=0}^l(-1)^j\Hilb_{N_j}(t).
\end{equation}
All factors $1-t^{w_i}$ are nonzero at $\zeta$.

Every finite graded $R/I$-module $N$ splits into its degree classes modulo
$M$; each class, regraded so that degree $M$ becomes degree one, has support
of dimension at most $\delta$, so $\Hilb_N$ has a pole of order at most
$\delta$ at $\zeta$. If the localization $N_\eta$ at the generic point $\eta$
of $Y$ is a $\C(Y)$-vector space, then
\begin{equation}\label{eq:leading}
 \lim_{t\to\zeta}(1-t^M)^\delta\Hilb_N(t)=e_Y\Tr(\zeta\mid N_\eta).
\end{equation}
Indeed, take a graded prime filtration of each degree class of $N$. Factors
supported in dimension less than $\delta$ contribute zero to the limit. Since
$Y$ is the unique component of dimension $\delta$, the remaining factors are
shifted copies of $\C[Y]$; one generated in degree $a$ contributes
$t^ah_Y(t^M)$, with leading coefficient $\zeta^ae_Y$. After localization
these factors give a filtration of $N_\eta$ by lines of character $\zeta^a$,
which gives the trace. That $e_Y>0$ is the positivity of the leading Hilbert
coefficient of the positively graded domain $\C[Y]$ regraded by $M$.

At a smooth closed point of $X$ on $Y$, diagonalize the finite-order action of
$\zeta$ on the cotangent space and lift eigenvectors to the maximal ideal by
the averaging projectors. They are formal coordinates, the completed local
ring is the power series ring in them with diagonal action, and the fixed
scheme is cut out by the coordinates with eigenvalue different from one. So
the fixed scheme is smooth and regularly embedded there, and its conormal
representation has no eigenvalue one. This applies at the generic point of
$Y$.

Put $K=\C(Y)$. In the regular local ring $R_\eta$, the ideal $I_\eta$ is the
maximal ideal, generated by a regular sequence of length $c$. There is an
exact sequence of $K$-representations of $\langle\zeta\rangle$
\[
 0\to L\to E=\bigoplus_{i=1}^lKe_i\to V=I_\eta/I_\eta^2\to0,\qquad
 e_i\mapsto f_i\bmod I_\eta^2 ,
\]
where $E$ has the characters $\zeta^{w_i}\ne1$ and $V$ the characters
$\nu_1,\dots,\nu_c\ne1$. Choosing among the $f_i$ a minimal generating set of
$I_\eta$, which is a regular sequence, and subtracting combinations of it from
the other generators changes the Koszul complex into the Koszul complex of a
regular sequence tensored with the Koszul complex on $l-c$ zeros; hence
$(N_j)_\eta\cong\bigwedge^j_KL$, naturally and therefore equivariantly. These
modules are vector spaces, so \eqref{eq:leading} applies, and
$\sum_j(-1)^j\Tr(\zeta\mid(N_j)_\eta)=\det(1-\zeta\mid L)$. Multiplying
\eqref{eq:euler} by $(1-t^M)^\delta$ and letting $t\to\zeta$,
\[
 \lim_{t\to\zeta}(1-t^M)^\delta\Hilb_R(t)
 =e_Y\frac{\det(1-\zeta\mid L)}{\prod_i(1-\zeta^{w_i})}
 =\frac{e_Y}{\det(1-\zeta\mid V)} ,
\]
which is \eqref{eq:residue}. The pole order is at most $\delta$ by the module
argument and exactly $\delta$ since the limit is not zero. Finally, if no
$a_i$ were divisible by $M$, the expression $P/\prod_i(1-t^{a_i})^{b_i}$
would be regular at $\zeta$.
\end{proof}

The Koszul lemmas used (null-homotopy, change of generators, vanishing of
higher Koszul homology for a regular sequence) are
\cite[Tags 0621, 062F]{Stacks}. Nonreduced structure of $X^\zeta$ at the
origin does no harm: its contributions are finite-dimensional.

\begin{lemma}\label{lem:orbitclosure}
Let $p\in X$ be a closed point other than the origin, and suppose that its
stabilizer in the grading group $\Gm$ is exactly $\mu_M$. Let $Y$ be the
closure of its $\Gm$-orbit, with reduced structure. Then $\C[Y]$ is
concentrated in degrees divisible by $M$, the set $\{k:\C[Y]_{Mk}\ne0\}$ is a
numerical semigroup (cofinite in $\Z_{\ge0}$), each nonzero
$\C[Y]_{Mk}$ is a line, and $e_Y=1$ with $\delta=1$.
\end{lemma}

\begin{proof}
Restriction to the orbit is the homomorphism $R\to\C[u]$, $f\mapsto f(p)u^d$
for $f\in R_d$; its image is $\C[Y]$, spanned by the $u^d$ with $d\in S=\{d:
f(p)\ne0\text{ for some }f\in R_d\}$, a semigroup since products of functions
nonzero at $p$ are nonzero at $p$. An element $\lambda\in\Gm$ fixes $p$ if and
only if $f(\lambda\cdot p)=\lambda^df(p)=f(p)$ for all homogeneous $f$, since
homogeneous functions separate the points of $X$; that is, if and only if
$\lambda^d=1$ for all $d\in S$. So the stabilizer is $\mu_{\gcd S}$, and
$\gcd S=M$. A subsemigroup of $M\Z_{\ge0}$ whose elements have greatest
common divisor $M$ contains all sufficiently large multiples of $M$. Hence
$h_Y(z)=1/(1-z)-P(z)$ for a polynomial $P$, and $e_Y=1$.
\end{proof}

\begin{lemma}\label{lem:lift}
Let $GL_n$ act linearly on a vector space $U$, let $R=\C[U]^{SL_n}$ be graded
so that $R_0=\C$, and let $\gamma$ be an odd integer such that every
$B\in GL_n$ acts on $X=\Spec R$ as the grading element $(\det B)^\gamma$. Let
$p\in X$ be a closed point other than the origin, fixed by a grading element
$\eta$ of order $2^a$, and let $y\in U$ lie in the closed $SL_n$-orbit over
$p$. Then there is a diagonalizable $A\in GL_n$ of finite two-power order with
$A\cdot y=y$ and $\ord(\det A)=2^a$.
\end{lemma}

\begin{proof}
As $\gamma$ is odd, there is $\theta$ of order $2^a$ with $\theta^\gamma=\eta$.
Choose $B\in GL_n$ with $\det B=\theta$. Then $B\cdot y$ lies in a closed orbit
over $\eta\cdot p=p$, and each fibre of the quotient map contains a unique
closed orbit \cite[Section I.1, p.~84]{Luna}; so $B\cdot y=S\cdot y$ with
$S\in SL_n$, and $C=S^{-1}B$ fixes $y$ with $\det C=\theta$. The semisimple
part $C_s$ of $C$ also fixes $y$ (an eigenvector of $C_s$ for an eigenvalue
$\ne1$ cannot be fixed by $C=C_sC_u$, whose unipotent part preserves the
eigenspaces of $C_s$) and has the same determinant. Diagonalize $C_s$ and let
$\Gamma\subset\C^*$ be the group generated by its eigenvalues, a finitely
generated abelian group, the product of its finite cyclic torsion subgroup and
a free group. The projection of $\Gamma$ onto the two-primary part of its
torsion subgroup, trivial on the free factor, is a homomorphism
$\Gamma\to\C^*$; applied to every eigenvalue it preserves every multiplicative
relation among them, in particular the relations expressing $C_s\cdot y=y$
weight by weight, and it fixes $\theta$. The resulting diagonal matrix $A$
fixes $y$, has finite two-power order, and $\det A=\theta$.
\end{proof}

\begin{lemma}\label{lem:cyclotomic}
Let $p$ be a prime and $d=p^re$ with $p\nmid e$. In $\F_p[t]$,
$\overline{\Phi_d}=\overline{\Phi_e}$ if $r=0$ and
$\overline{\Phi_d}=\overline{\Phi_e}^{\,p^{r-1}(p-1)}$ if $r\ge1$. The roots of
$\overline{\Phi_e}$ in $\overline\F_p$ are the elements of multiplicative order
$e$, each simple.
\end{lemma}

\begin{proof}
For $r\ge1$, $\Phi_{ep^r}(t)\Phi_e(t^{p^{r-1}})=\Phi_e(t^{p^r})$ in $\Z[t]$
(group the roots on the right by the $p$-part of their order). Modulo $p$,
$t\mapsto t^{p^j}$ becomes the $p^j$-th power; cancel
$\overline{\Phi_e}^{\,p^{r-1}}$. The factorization $t^e-1=\prod_{b\mid e}
\Phi_b(t)$ persists modulo $p$, $t^e-1$ is separable there, and induction on
$e$ assigns to $\overline{\Phi_e}$ exactly the roots of order $e$.
\end{proof}

\begin{theorem}\label{thm:transfer}
Let $F,G\in\Q(t)\cap\Z[[t]]$ have all their poles at roots of unity, let $p$ be
a prime, $p\nmid M$, and suppose $F\equiv G\pmod{p\Z[[t]]}$. Suppose $G$ has a
pole of order $\delta\ge1$ at a primitive $M$-th root of unity $\zeta$, no pole
at any primitive $Mp^r$-th root of unity for $r\ge1$, and that
$C=\lim_{t\to\zeta}(1-t^M)^\delta G(t)$ is a unit in $\Z[\zeta]_{\mathfrak p}$
for some prime $\mathfrak p$ over $p$. Then $F$ has a pole at a primitive
$Mp^r$-th root of unity for some $r\ge0$, hence at all of them; and every
expression $F(t)=P(t)/\prod_i(1-t^{a_i})^{b_i}$ with $P\in\Q[t]$ and $a_i,b_i>0$
has some $a_i$ divisible by $Mp^r$, in particular $a_i\ge M$.
\end{theorem}

\begin{proof}
A reduced expression of $F$ has a denominator that is, up to sign, a product
of cyclotomic polynomials; normalize its constant term to $1$. The numerator is
then integral, being the product of the integral series $F$ with an integral
polynomial. So $F=N_F/D_F$ and $G=N_G/D_G$ with $N_\bullet,D_\bullet\in\Z[t]$,
$D_\bullet(0)=1$, coprime over $\Q$, the $D_\bullet$ signed products of
cyclotomic polynomials. Write $D_G=\varepsilon\Phi_M^\delta D_0$. By
Lemma~\ref{lem:cyclotomic} and the hypothesis on poles of order $Mp^r$, the
reduction $\overline{D_0}$ does not vanish at $\alpha=\zeta\bmod\mathfrak p$,
which has exact order $M$. With $U_M=(1-t^M)/\Phi_M$, separability of
$t^M-1$ modulo $p$ gives $\overline{U_M}(\alpha)\ne0$. So $D_0(\zeta)$ and
$U_M(\zeta)$ are $\mathfrak p$-units, and
$C=U_M(\zeta)^\delta N_G(\zeta)/(\varepsilon D_0(\zeta))$ forces
$N_G(\zeta)$ to be a unit: $\overline{N_G}(\alpha)\ne0$ while
$\overline{D_G}$ vanishes at $\alpha$. So $\overline G\in\F_p(t)$ has a pole
at $\alpha$. Reduction of both fractions is defined in $\F_p(t)$ and in
$\F_p[[t]]$, and the map from such fractions to series is injective; the
congruence gives $\overline F=\overline G$ in $\F_p(t)$, so
$\overline{D_F}(\alpha)=0$. Some cyclotomic factor $\Phi_d$ of $D_F$ has
$\overline{\Phi_d}(\alpha)=0$, and by Lemma~\ref{lem:cyclotomic} the prime-to-$p$
part of $d$ is $M$: $d=Mp^r$. As $N_F,D_F$ are coprime over $\Q$, $F$ has a
pole at a primitive $d$-th root of unity, and by rationality over $\Q$ at all
of them. A denominator $\prod_i(1-t^{a_i})^{b_i}$ must vanish there, so
$d\mid a_i$ for some $i$.
\end{proof}

\begin{remark}
The hypothesis on the poles of order $Mp^r$ cannot be dropped: $G(t)=1+pt^{pM}/
(1-t^{pM})$ has integral nonnegative coefficients and $\lim_{t\to\zeta}
(1-t^M)G(t)=1$, but $\overline G=1$. Nor can the conclusion be improved to
$r=0$: for $p=3$, $M=2^a\ge4$ and $z=t^{M/2}$, the series
$F=(1+z)/(1-z+z^2)+2t/(1-t)^2$ and $G=1/(1+z)+2t/(1-t)^2$ have nonnegative
integral coefficients, $F-G=3z/(1+z^3)$, the residue of $G$ at the primitive
$M$-th roots is $2$, a $3$-adic unit, but $F$ has its poles there only at order
$3M$, since $1-z+z^2=\Phi_{3M}(t)$.
\end{remark}

\section{Cubic supports fixed by a two-power diagonal matrix}\label{sec:support}

Let $F\in\C[x_1,\dots,x_n]$ be a cubic form and $B$ the integer matrix whose
columns are the exponent vectors of its monomials with nonzero coefficient
(its \emph{support}). We say $F$ satisfies the \emph{span condition} if
$\one=(1,\dots,1)^t$ lies in the rational column span of $B$.

\begin{lemma}\label{lem:span}
Every cubic form at which some $SL_n$-invariant of positive degree does not
vanish satisfies the span condition, in every coordinate basis.
\end{lemma}

\begin{proof}
Some coefficient monomial of the invariant is nonzero at $F$ and so involves
only coefficients of supported monomials. Invariance under the diagonal torus
of $SL_n$ makes the sum of the corresponding exponent vectors, counted with
multiplicity, a vector with all coordinates equal, and positive since the
degree is positive.
\end{proof}

If a diagonal matrix $A$ with eigenvalues $\alpha_1,\dots,\alpha_n$ of finite
two-power order fixes $F$, the character $\psi:\Z^n\to\C^*$,
$u\mapsto\prod_i\alpha_i^{u_i}$, kills the lattice generated by the support,
and $\psi(\one)=\det A$. If $K$ is a set of support columns forming a basis of
the rational support span, the equation $Kz=\one$ has a unique rational
solution; by Cramer's rule applied to a nonzero maximal minor $K_I$, the least
$d\ge1$ with $dz$ integral divides $\det K_I$, and $d\one=K(dz)$ is in the
support lattice. Hence
\begin{equation}\label{eq:cramer}
 \ord(\det A)\le2^{v_2(\det K_I)} .
\end{equation}

\paragraph{The support graph.}
A monomial $x_i^2x_j$ ($i\ne j$) gives an edge between $i$ and $j$ with
coefficient two at $i$ and one at $j$; a cube $x_i^3$ gives a loop with
coefficient three; all $n$ vertices are included. Every other monomial is of
the form $x_ix_jx_k$ with three distinct indices (a \emph{triple}). For a
spanning tree of a connected component with $v$ vertices, the $v\times(v-1)$
matrix $T$ of its columns has rank $v-1$ (remove a leaf: its row has a single
nonzero entry, one or two). So a graph component has rank $v$ or $v-1$ (an
isolated vertex has rank one if its cube occurs, zero otherwise). Deleting a
row $r$ of $T$ and removing leaves other than $r$ one at a time shows
\begin{equation}\label{eq:treeminor}
 \det T^{(r)}=\pm2^{e_r},\qquad 0\le e_r\le v-1 .
\end{equation}
The signed cofactor vector $u$ of $T$ spans the kernel of $T^t$; along an edge
with coefficients two and one it satisfies $2u_i+u_j=0$, so adjacent entries
differ in absolute value by a factor two and have opposite signs:
\begin{equation}\label{eq:cofactor}
 u=\varepsilon2^{e}\bigl((-2)^{h_v}\bigr)_v,\qquad
 \varepsilon=\pm1,\quad e\ge0,\quad h_v\ge0,\quad\min_vh_v=0 .
\end{equation}

\begin{proposition}\label{prop:graphsupport}
Suppose the support of $F$ consists of cubes and monomials $x_i^2x_j$ only, and
$F$ satisfies the span condition. If a diagonal matrix $A$ of finite two-power
order fixes $F$, then $\ord(\det A)\le2^{n-1}$. If equality holds, then after
a permutation of the variables the support is exactly
\begin{equation}\label{eq:chainsupport}
 x_1^2x_2,\ x_2^2x_3,\ \dots,\ x_{n-1}^2x_n,\ x_n^3,
\end{equation}
the eigenvalues of $A$ are $\xi,\xi^{-2},\xi^4,\dots,\xi^{(-2)^{n-1}}=1$ with
$\xi$ of order $2^{n-1}$, and a diagonal substitution turns $F$ into the chain
cubic $F_n=x_1^2x_2+\dots+x_{n-1}^2x_n+x_n^3$.
\end{proposition}

\begin{proof}
The support lattice, the span condition and $\psi$ split over the connected
components of the support graph, and $\det A$ is the product of the values of
$\psi$ on the all-ones vectors $\one_v$ of the components. Each value has
two-power order; we show it is at most $2^{v-1}$ for a component with $v$
vertices, and at most $2^{v-2}$ if $v\ge2$ and the component is rank-deficient.
An isolated vertex must carry a cube (span condition), and $3e_i$ in the
lattice makes $\psi(e_i)$ trivial.

\emph{Rank $v-1$.} The span of the component is that of a spanning tree $T$,
and $Tc=\one_v$ has a unique rational solution. By Cramer's rule for every
$T^{(r)}$, the denominator $q$ of $c$ divides every $2^{e_r}$, and
$q\one_v=T(qc)$ lies in the lattice. By \eqref{eq:cofactor} adjacent $e_r$
differ by one; all are at most $v-1$, so their minimum is at most $v-2$, and
$q\le2^{v-2}$.

\emph{Rank $v$.} One further column extends $T$ to a nonsingular square
matrix $C$, whose graph is connected and unicyclic (the cycle may be a loop or
a double edge). The order of $\psi(\one_v)$ divides $|\det C|$, so it is at
most $2^{v_2(\det C)}$. Expanding along leaf rows contributes a factor one or
two for each of the $v-l$ tree edges off the cycle of length $l$. A loop
contributes the odd factor three, so $v_2(\det C)\le v-1$. A cycle of length
$l\ge2$ has determinant $\pm(2^a+(-1)^{l-1}2^{l-a})$ for some $0\le a\le l$;
it is nonzero, so $a\ne l-a$ when $l$ is even, and its valuation is
$\min(a,l-a)\le\lfloor(l-1)/2\rfloor$. Then $v_2(\det C)\le v-l+
\lfloor(l-1)/2\rfloor\le v-2$.

The product of two-power orders has order at most the largest one, and every
component has at most $n$ vertices; this gives $\ord(\det A)\le2^{n-1}$. In the
case of equality, some component has $n$ vertices, so the graph is connected,
of full rank, and the chosen unicyclic basis has a loop, at a vertex we call
the root; and every leaf-expansion factor is two. Rooting the tree at the
loop, each tree edge has coefficient two at its child and one at its parent,
and the eigenvalue relations are $\alpha_{\rm root}^3=1$ and
$\alpha_{\rm parent}=\alpha_{\rm child}^{-2}$. As the eigenvalues have two-power
order, $\alpha_{\rm root}=1$, and a vertex at depth $d$ has eigenvalue of order
at most $2^d$; if the depth of the tree is $h$, then $\det A$ has order at most
$2^h$. So $h\ge n-1$, and the tree is a path rooted at one end. The deepest
eigenvalue $\xi$ has order $2^{n-1}$, and the eigenvalues along the path are
$\xi,\xi^{-2},\dots,\xi^{(-2)^{n-1}}=1$. In particular $A$ has order $2^{n-1}$.

These eigenvalues fix no other cubic monomial. Indeed, the exponents
$(-2)^i\bmod2^{n-1}$, $0\le i\le n-2$, have distinct two-adic valuations and
the last exponent is $0$. In a sum of three of them that vanishes modulo
$2^{n-1}$, the least valuation among the nonzero ones cannot occur exactly
once, nor three times (three equal nonzero terms have the valuation of one of
them, since three is odd); so it occurs twice, at one index $i$, and the third
term is $-2(-2)^i$, the next chain exponent, or zero if $i=n-2$. Three zero
terms give $x_n^3$. So the monomials fixed by $A$ are exactly
\eqref{eq:chainsupport}, and they all occur, being the columns of the basis.
Rescale $x_n$ by a cube root and then each $x_i$, $i=n-1,\dots,1$, by a square
root to make all coefficients one.
\end{proof}

\begin{theorem}\label{thm:fullcubicsupport}
Let $F$ be a cubic form in $n\ge2$ variables satisfying the span condition, and
let a diagonal matrix $A$ of finite two-power order fix $F$. Then
$\ord(\det A)\le2^{n-1}$. If equality holds, then after a permutation of the
variables the support of $F$ is exactly \eqref{eq:chainsupport}, $A$ has order
$2^{n-1}$, and a diagonal substitution turns $F$ into $F_n$. The bound is
attained by $F_n$.
\end{theorem}

\begin{proof}
Let $G$ be the matrix of the graph columns (cubes and $x_i^2x_j$). By the
spanning-tree argument, if the graph has $f$ components of full rank and $c$
rank-deficient components, with $v_i$ vertices, then $\rank G=n-c$. In each
full component choose a square basis $C_i$ (a spanning tree and one more graph
column); the proof of Proposition~\ref{prop:graphsupport} gives
\begin{equation}\label{eq:fullcomp}
 v_2(\det C_i)\le v_i-1 .
\end{equation}
In each deficient component choose a spanning tree $T_i$; its signed cofactor
vector is $2^{e_i}u_i$ with $(u_i)_v=\varepsilon_i(-2)^{h_{iv}}$ as in
\eqref{eq:cofactor} (for an isolated vertex without cube, $e_i=h_{iv}=0$).
Put $b_i=v_i-1-e_i$ and $H_i=\max_vh_{iv}$. Since a tree minor is a product of
$v_i-1$ factors one or two, $0\le H_i\le b_i$; and if $H_i=b_i$, the vertex of
height $H_i$ is unique: rooting the tree there, its deleted-row minor has
absolute value $2^{v_i-1}$, so every factor is two, each edge has coefficient
two at its child and one at its parent, $h_{\rm parent}=h_{\rm child}+1$, and
heights decrease strictly away from the root.

Normalize the entries of $u_i$ by $2^{b_i}$. For sums over distinct vertices
$E$ of one deficient component,
\begin{equation}\label{eq:smallsums}
 \Bigl|\sum_{v\in E}\frac{(u_i)_v}{2^{b_i}}\Bigr|\le1,\ \tfrac54,\ \tfrac32
 \qquad\text{for }|E|=1,2,3 .
\end{equation}
If $H_i<b_i$, all entries have absolute value at most $1/2$. If $H_i=b_i$,
there is one entry of absolute value one, which we make positive; the entries
of absolute value $1/2$ are at height $H_i-1$ and so negative; all others have
absolute value at most $1/4$. A pair avoiding the top entry has sum at most one
in absolute value, a pair containing it at most $5/4$; a triple avoiding it at
most $3/2$; a triple containing it gives $0$ if both others are $-1/2$, at
most $3/4$ if one is, and at most $1+1/4+1/4$ otherwise.

The bases $C_i$ and trees $T_i$ form a rational basis of the graph columns.
Extend it by $q$ triple columns $\beta_1,\dots,\beta_q$ to a basis of all
support columns; then $0\le q\le c$ and the support rank is $r=n-c+q$. If
$q=0$, the graph part satisfies the span condition by itself, and
Proposition~\ref{prop:graphsupport} applies: it gives the bound, and in case
of equality the chain eigenvalues, which fix no triple monomial; so $F$ has no
triple monomials and the theorem holds. Assume $q>0$.

Define the $c\times q$ matrix $D_{ij}=u_i^t(\beta_j|_i)$, where $\beta_j|_i$ is
the restriction to the vertices of deficient component $i$. The map sending a
vector to its $c$ pairings with the $u_i$ is surjective with kernel the
rational span of the graph columns (on a full component the graph columns span
all coordinates; on a deficient one $\ker u_i^t$ is the tree span), so $D$ has
rank $q$; choose $q$ deficient components $J$ with $\det D_J\ne0$. Let $K$ be
the $n\times r$ basis matrix. Keep all rows of full components and of the
components in $J$; in every other deficient component delete one row at a
vertex of height zero. The resulting $r\times r$ minor satisfies
\begin{equation}\label{eq:minorfactor}
 \det K_I=\pm\Bigl(\prod_{\text{full }i}\det C_i\Bigr)
 2^{\sum_{\text{deficient }i}e_i}\det D_J .
\end{equation}
Indeed, in a full component eliminate the triple entries with its graph
columns and factor out $\det C_i$; in a deficient component outside $J$ the
retained tree block is square with determinant $\pm2^{e_i}$ (the deleted row
had height zero); eliminate and factor out. These operations stay inside
components. What remains are the tall tree blocks of $J$ and the $q$ triple
columns; expanding along maximal minors of each tree block leaves one row per
component, given by the cofactor vector $2^{e_i}u_i$, and the remaining
$q\times q$ determinant is $\det D_J$.

Divide row $i$ of $D_J$ by $2^{b_i}$. Each column then involves at most three
distinct vertices; by \eqref{eq:smallsums} its squared Euclidean norm is at
most $3$ (three components), $(5/4)^2+1=41/16$ (two and one) or $(3/2)^2$ (all
in one). Hadamard's inequality gives
\begin{equation}\label{eq:hadamard}
 0<|\det D_J|\le2^{\sum_{i\in J}b_i}3^{q/2} .
\end{equation}
If $q<c$, or $q=c$ and $f>0$, then by \eqref{eq:fullcomp},
\eqref{eq:minorfactor} and \eqref{eq:hadamard},
\[
 v_2(\det K_I)\le\sum_{\rm full}(v_i-1)+\sum_{\rm def}e_i+\sum_{i\in J}b_i
 +\tfrac q2\log_23
 =n-f-c-\!\!\sum_{\substack{{\rm def}\ i\notin J}}\!\!b_i+\tfrac q2\log_23<n-f-c+q,
\]
strictly since $q>0$ and $\sqrt3<2$; the right side is at most $n-1$, so
$v_2(\det K_I)\le n-2$. If $q=c$ and $f=0$, then $K$ is square and every
column has coordinate sum three, so $\one^tK=3\one^t$ and $3\mid\det K$, hence
$3\mid\det D$; with $\sum_i(e_i+b_i)=n-c$,
\[
 2^{v_2(\det K)}\le2^{\sum e_i}\frac{|\det D|}3\le2^{n-c}3^{c/2-1}<\frac{2^n}3
 <2^{n-1}.
\]
In both cases \eqref{eq:cramer} gives $\ord(\det A)\le2^{n-2}$. So equality
occurs only with $q=0$. Finally $F_n$ is fixed by
$\diag(\xi,\xi^{-2},\dots,\xi^{(-2)^{n-1}})$ with $\xi$ of order $2^{n-1}$,
whose determinant $\xi^{S}$, $S=\sum_{i<n}(-2)^i$, has order $2^{n-1}$ since
$S$ is odd.
\end{proof}

\section{Kronecker coefficients}\label{sec:kronecker}

Let $W=\C^n$ with basis $e_1,\dots,e_n$, and let $R_n=\C[W\otimes W\otimes W]
^{SL_n^3}$, graded so that $(R_n)_s$ is the space of invariants of degree $ns$
(invariants occur only in degrees divisible by $n$). By Schur--Weyl duality
$\dim(R_n)_s=g((s^n),(s^n),(s^n))=k_n(s)$, so $K_n=\Hilb_{R_n}$. Let $\tau$ be
the cyclic permutation of the three tensor factors; it acts on $R_n$ as a
graded algebra automorphism of order three.

Let $\mathcal U_n=W\oplus\mathbb S_{(2,1)}W$, give its coordinate ring the
grading in which linear functions on $W$ have weight one and those on
$\mathbb S_{(2,1)}W$ weight three, and let $b_n(s)$ be the dimension of the
$SL_n$-invariants of weighted degree $ns$ and $B_n(t)=\sum_sb_n(s)t^s$. (All
invariants have weighted degree divisible by $n$, as the scalars $\mu_n$ act
on weighted degree $D$ by $\zeta^D$.) Then $B_n=\Hilb_{T_n}$ with
$T_n=\C[\mathcal U_n]^{SL_n}$ graded by $s$.

\begin{theorem}\label{thm:K023}
For every $n\ge2$ and $s\ge0$,
\[
 \tr\bigl(\tau\mid(R_n)_s\bigr)=b_n(s),\qquad k_n(s)\equiv b_n(s)\pmod3 ,
\]
and $k_n(s)-b_n(s)$ is a nonnegative multiple of three. Hence
$K_n(t)\equiv B_n(t)\pmod{3\Z[[t]]}$.
\end{theorem}

\begin{proof}
We may compute with polynomial representations of $W$ instead of $W^*$; the
simultaneous contragredient preserves the dimensions of invariant spaces and
commutes with permutations of the factors. In degree $d=ns$, Schur--Weyl
duality gives, compatibly with $\tau$,
\[
 (R_n)_s\simeq\bigl([\lambda]\otimes[\lambda]\otimes[\lambda]\bigr)^{S_d},
 \qquad\lambda=(s^n),
\]
since $\Sym^d(W\otimes W\otimes W)$ is the diagonal $S_d$-invariant part of
$(W^{\otimes d})^{\otimes3}$ and $(W^{\otimes d})^{SL_n}$ is the Specht module
$[\lambda]$. For an endomorphism $A$ of a vector space, the trace of the cyclic
permutation composed with $A^{\otimes3}$ is $\tr(A^3)$ (contract the cyclic
indices in a basis). Averaging over $S_d$,
\[
 \tr\bigl(\tau\mid(R_n)_s\bigr)=\frac1{d!}\sum_{\pi\in S_d}\chi_\lambda(\pi^3).
\]
By the Frobenius formula $p_\mu=\sum_\lambda\chi_\lambda(\mu)s_\lambda$, this
average is the coefficient of $s_\lambda z^d$ in
$\sum_d\frac{z^d}{d!}\sum_{\pi\in S_d}p_{{\rm type}(\pi^3)}$. Cubing a
permutation keeps a $j$-cycle with $3\nmid j$ and splits a $3r$-cycle into three
$r$-cycles, so by the exponential formula the series is
\[
 \exp\Bigl(\sum_{3\nmid j}\frac{z^jp_j}j+\sum_{r\ge1}\frac{z^{3r}p_r^3}{3r}\Bigr)
 =\exp\Bigl(\sum_{r\ge1}\frac{z^rp_r}r+\sum_{r\ge1}\frac{z^{3r}}r\,
 p_r[s_{(2,1)}]\Bigr),
\]
using $p_r[s_{(2,1)}]=(p_r^3-p_{3r})/3$, which follows from
$s_{(2,1)}=(p_1^3-p_3)/3$. The right side is the character of the symmetric
algebra on $W\oplus\mathbb S_{(2,1)}W$ with weights one and three
\cite[Chapter I, \S\S7--8]{Macdonald}; specializing to $n$ variables and taking
the coefficient of $s_{(s^n)}$ gives $b_n(s)$. Finally, $\tau$ is defined over
$\Q$ on $(R_n)_s$, and an operator of order three on a rational vector space has
characteristic polynomial $(X-1)^a(X^2+X+1)^b$, dimension $a+2b$ and trace
$a-b$; the difference is $3b\ge0$.
\end{proof}

We realize $\mathbb S_{(2,1)}W$ as the degree-three part of the free Lie
algebra on $W$ inside $W^{\otimes3}$, with $[a,b]=ab-ba$: the map
$W\otimes\Lambda^2W\to W^{\otimes3}$, $a\otimes(b\wedge c)\mapsto[a,[b,c]]$,
has kernel $\Lambda^3W$ by the Jacobi identity, and its image is
$\mathbb S_{(2,1)}W$. Its weights are $2e_i+e_j$ ($i\ne j$), with multiplicity
one, and $e_i+e_j+e_k$ ($i,j,k$ distinct), with multiplicity two. Put
\[
 z_n=(e_n,F_n)\in\mathcal U_n,\qquad F_n=\sum_{i=1}^{n-1}[e_i,[e_i,e_{i+1}]],
 \qquad M=2^{n-1}.
\]

\begin{lemma}\label{lem:liezero}
For every $n\ge2$, the stabilizer of $z_n$ in $\mathfrak{gl}_n$ is zero, so the
stabilizer of $z_n$ in $GL_n$ is finite.
\end{lemma}

\begin{proof}
Write $Xe_j=\sum_ax_{aj}e_a$ and suppose $Xz_n=0$; then $Xe_n=0$, and $X$ acts
on $F_n$ as a derivation of the brackets. For $r\ne s$, the weight $2e_r+e_s$
is spanned by $L_{rs}=[e_r,[e_r,e_s]]$, and its coefficient in $XF_n$ is
\[
 \mathbf1_{r<n}x_{s,r+1}+2\cdot\mathbf1_{s=r+1}x_{rr}-\mathbf1_{s=r-1}x_{r,r-1},
\]
the last term coming from $[e_r,[e_{r-1},e_r]]=-L_{r,r-1}$. Taking $s=r+1$ gives
$2x_{rr}+x_{r+1,r+1}=0$, and $x_{nn}=0$, so the diagonal vanishes. For column
$j\ge2$ the only possibly nonzero entries are then $x_{j-1,j}$ and
$x_{j-2,j}=x_{j-1,j-2}$ ($j\ge3$); for $j\ge4$ the latter is an entry below the
diagonal in column $j-2\ge2$, hence zero. So besides the first column and the
superdiagonal only $x_{13}=x_{21}$ ($n\ge3$) can be nonzero. Put
$y_j=x_{j-1,j}$; $Xe_n=0$ gives $y_n=0$. If $n=2$, $Xe_2=0$ also kills $x_{12}$,
and the coefficient of $L_{21}$ gives $x_{21}=0$. Let $n\ge3$.

For $a\ge4$, the weight $e_1+e_2+e_a$ receives only the term of index one,
with coefficient $x_{a1}$ times the nonzero vector
$[e_a,[e_1,e_2]]+[e_1,[e_a,e_2]]$; so $x_{a1}=0$. For $3\le j\le n-1$, the
weight $e_{j-1}+e_j+e_{j+1}$ receives only the two adjacent terms, with
coefficients $y_j$ and $x_{j+1,j-1}=0$, the former multiplying the nonzero
vector $[e_{j-1},[e_j,e_{j+1}]]+[e_j,[e_{j-1},e_{j+1}]]$; so $y_j=0$. In the
weight $e_1+e_2+e_3$, with the basis $A=[e_1,[e_2,e_3]]$,
$B=[e_2,[e_1,e_3]]$ and $[e_3,[e_1,e_2]]=-A+B$, the contribution is
$y_2(A+B)+x_{31}(-2A+B)$; so $y_2-2x_{31}=y_2+x_{31}=0$ and
$x_{31}=y_2=0$. Finally $x_{13}=x_{21}$ vanishes by $Xe_3=0$ if $n=3$, and for
$n\ge4$ because the weight $e_1+e_3+e_4$ receives only the derivative of the
repeated $e_3$ in the term of index three, with coefficient $x_{13}$ times a
nonzero vector. So $X=0$. An algebraic group in characteristic zero has the
dimension of its Lie algebra.
\end{proof}

\begin{theorem}\label{thm:K024}
Let $n\ge2$, $M=2^{n-1}$, $\zeta$ a primitive $M$-th root of unity, and
$X_n=\Spec T_n$.
\begin{enumerate}
\item The $SL_n$-orbit of $z_n$ is closed; its stabilizer in $GL_n$ is cyclic of
order $M$, generated by
\begin{equation}\label{eq:chainelement}
 g=\diag(\zeta^{w_1},\dots,\zeta^{w_n}),\qquad w_i=(-2)^{i-1}\bmod M ,
\end{equation}
on which $\det$ is faithful; its stabilizer in $SL_n$ is trivial.
\item The nonzero closed points of $X_n$ fixed by the grading group $\mu_M$
form one orbit of the grading group $\Gm$, the orbit of the image $p_n$ of
$z_n$, whose stabilizer in $\Gm$ is exactly $\mu_M$. $X_n$ is smooth at $p_n$,
and the fixed scheme $X_n^{\zeta}$ is smooth of dimension one there.
\item $B_n$ has a simple pole at every primitive $M$-th root of unity $\xi$,
with $\lim_{t\to\xi}(1-t^M)B_n(t)$ a unit in $\Z[\xi,1/2]$.
\item If the two-part of an integer $d$ is at least $M$ and $d\ne M$, then
$B_n$ is regular at every primitive $d$-th root of unity; in particular at the
primitive $M3^r$-th roots for $r\ge1$.
\end{enumerate}
\end{theorem}

\begin{proof}
\emph{Closed orbit.} For $1\le i<n$ put $q_i=(1-(-1/2)^i)/3$ and $q_n=1-q_{n-1}$;
all are positive, and with $a_i=2e_i+e_{i+1}$ ($i<n$), $a_n=e_n$, one has
$\sum_iq_ia_i=\one$ ($2q_1=1$, $q_{i-1}+2q_i=1$, $q_{n-1}+q_n=1$). The words
$[e_i,[e_i,e_{i+1}]]$ have squared norm six in the Hermitian metric of
$W^{\otimes3}$; give the word $i$ the coefficient $\sqrt{q_i/6}$ and $e_n$ the
coefficient $\sqrt{q_n}$. The diagonal part of the moment map of $SU_n$ is then
$\one$ minus its trace part, which vanishes; its off-diagonal part vanishes
because no two of the weights $a_1,\dots,a_n$ in the same summand differ by a
root $e_a-e_b$, and the summands are orthogonal and preserved. So this vector
is a zero of the moment map, and its $SL_n$-orbit is closed by the
Kempf--Ness criterion \cite[Theorem 4, p.~1]{Losev}. The weights $a_i$ are
linearly independent (triangular, determinant $2^{n-1}$), so the diagonal torus
acts transitively on vectors with these weights and nonzero coefficients; it
maps our vector to $z_n$, and the orbit of $z_n$ is closed as well.

\emph{Stabilizer.} The weights $w_i$ are distinct modulo $M$, $w_n=0$ and
$2w_i+w_{i+1}\equiv0$, so $g$ fixes $z_n$, and $\det g=\zeta^{\sum w_i}$ with
$\sum_iw_i$ odd. Let $L=\Stab_{GL_n}(z_n)$, finite by Lemma~\ref{lem:liezero}.
As $g$ has simple spectrum, its centralizer in $L$ is diagonal; a diagonal
element of $L$ satisfies $d_n=1$ and $d_i^2d_{i+1}=1$, so the centralizer is
$\langle g\rangle$, of order $M$. Column orthogonality gives
$\sum_{\rho\in{\rm Irr}(L)}|\rho(g)|^2=M$; the powers $\det^k$, $0\le k<M$, are
distinct linear characters (as $\det g$ has order $M$) and exhaust this sum. So
every linear character of $L$ is a power of $\det$, and every nonlinear
irreducible character vanishes at $g$. Let $K=\ker(\det|_L)$. A representation
of $g$ has trace zero exactly when the multiplicities of $\eta$ and $-\eta$
agree for every eigenvalue $\eta$ (use the basis $1,\zeta,\dots,\zeta^{M/2-1}$
of $\Q(\zeta)$ and $\zeta^{M/2}=-1$). The spectrum \eqref{eq:chainelement} has
exactly one pair $\{\eta,-\eta\}$, on $P=\langle e_{n-1},e_n\rangle$, since
the $w_i$ with $i\le n-2$ have distinct two-adic valuations less than $n-2$.
So every nonlinear constituent of $W$ lies in $P$; the linear constituents are
powers of $\det$, on which $K$ acts trivially; $P$ is $L$-stable, $K$ acts on
it with determinant one and fixes $e_n$, and a finite-order element of
$SL_2(\C)$ fixing a nonzero vector is the identity. So $K=1$ and
$L=\langle g\rangle$; the $SL_n$-stabilizer is trivial.

\emph{The grading.} The scalar $\lambda I\in GL_n$ acts on $\mathcal U_n$ as the
weight grading, and on $(T_n)_s$ by $\lambda^{\pm ns}=(\det\lambda I)^{\pm s}$;
so each $B\in GL_n$ acts on $X_n$ as the grading element $(\det B)^{\pm1}$.
Hence the grading stabilizer of $p_n$ is $\det(L)=\mu_M$, and every point of
$\Gm\cdot p_n$ is fixed by $\mu_M$.

\emph{Fixed points.} Let $p\ne0$ be a closed point of $X_n$ fixed by a grading
element of order $2^a\ge M$ and $y$ in its closed orbit. By
Lemma~\ref{lem:lift} (with $\gamma=\pm1$) some diagonal $A$ of two-power order
fixes $y$ with $\ord(\det A)=2^a$. In an eigenbasis of $A$ consider the weights
of the nonzero components of $y$: linear weights $e_i$, and weights $2e_i+e_j$,
$e_i+e_j+e_k$ of $\mathbb S_{(2,1)}W$. Replace each linear weight $e_i$ by
$3e_i$; this preserves the rational span and, for eigenvalues of two-power
order, the condition of being fixed ($\alpha_i=1\iff\alpha_i^3=1$). The span
condition holds by the argument of Lemma~\ref{lem:span}, applied to a
positive-degree invariant nonzero at $y$. The resulting columns are the support
of a cubic form fixed by $A$, and Theorem~\ref{thm:fullcubicsupport} gives
$2^a\le M$; so $a=n-1$, and the support is the chain \eqref{eq:chainsupport}
after a permutation. The cube $3e_n$ comes from the linear coordinate $e_n$,
since $\mathbb S_{(2,1)}W$ has no weight $3e_i$, and the other columns from the
one-dimensional weight spaces of $2e_i+e_{i+1}$. So $y$ is diagonally
equivalent to $z_n$, and $p\in\Gm\cdot p_n$.

\emph{Smoothness and the fixed tangent space.} Since the closed orbit of $z_n$
has trivial stabilizer in $SL_n$, Luna's slice theorem
\cite[Section III.1, Th\'eor\`eme du slice \'etale and Remarque 1$^\circ$, pp.~97--98]{Luna} shows that $X_n$ is smooth at $p_n$ with
tangent space $\mathcal U_n/\mathfrak{sl}_n\cdot z_n$, equivariantly for $g$.
The fixed space of $g$ on $\mathcal U_n$ has dimension $n$. Recall
$v_2(w_i)=i-1$ for $i<n$ and $w_n=0$. A linear weight $e_i$ is fixed only for
$i=n$. A weight $2e_i+e_j$ is fixed only for $j=i+1$: if $j=n$, then
$2w_i\equiv0$ forces $i=n-1$; if $i=n$, then $w_j\not\equiv0$; if $i,j<n$ and
$j\ne i+1$, the valuations $i$ and $j-1$ of $2w_i$ and $w_j$ differ, so the sum
has valuation at most $n-2$; and $2w_i+w_{i+1}=0$. A weight $e_i+e_j+e_k$ is
never fixed, since among its nonzero terms the least valuation occurs once.
As $g$
has simple spectrum, $\dim\mathfrak{sl}_n^g=n-1$. Taking fixed vectors is exact
for a finite group in characteristic zero, so $\dim(T_{p_n}X_n)^g=1$, and as in
the proof of Proposition~\ref{prop:fixedcomponentpole} the fixed scheme is
smooth of dimension one at $p_n$. Translating by the grading group, the same
holds along $\Gm\cdot p_n$.

\emph{The pole.} Let $Y$ be the closure of $\Gm\cdot p_n$. Its only other point
is the origin, since the grading is positive; so the support of
$X_n^{\zeta}$ is $Y$, a curve, and $X_n$ is smooth at its generic point. By
Lemma~\ref{lem:orbitclosure}, $e_Y=1$. Proposition~\ref{prop:fixedcomponentpole}
gives a simple pole at $\zeta$ with $\lim(1-t^M)B_n=1/\prod_j(1-\nu_j)$, where
the $\nu_j$ are the eigenvalues of $\zeta$ on the conormal space, all
different from one and all $M$-th roots of unity. For an $M$-th root of unity
$\nu\ne1$ of order $2^b$, $(1-\nu)\sum_{j<2^{b-1}}\nu^j=1-\nu^{2^{b-1}}=2$, so
$1-\nu$ is a unit in $\Z[\zeta,1/2]$. The same holds at every primitive
$M$-th root by Galois conjugation, $B_n$ having rational coefficients.

\emph{No other poles of large two-part.} If the two-part of $d$ is at least $M$
and $d\ne M$, a nonzero point fixed by a primitive $d$-th root would be fixed by
its two-primary part, of order at least $M$, hence lie on $\Gm\cdot p_n$, whose
grading stabilizer is $\mu_M\not\ni$ the primitive $d$-th roots. So the fixed
locus is the origin, and Lemma~\ref{lem:isolated} applies.
\end{proof}

\begin{theorem}\label{thm:kronecker}
For every $n\ge2$ there is $r\ge0$ such that $K_n(t)=\sum_sg((s^n),(s^n),
(s^n))t^s$ has a pole at every primitive root of unity of order $2^{n-1}3^r$.
Every expression $K_n(t)=P(t)/\prod_i(1-t^{a_i})^{b_i}$ with $P\in\Q[t]$ and
positive integers $a_i,b_i$ has some $a_i$ divisible by $2^{n-1}3^r$; in
particular every positive form of $K_n$ has modular index at least $2^{n-1}$.
\end{theorem}

\begin{proof}
Both $K_n$ and $B_n$ are Hilbert series of finitely generated positively graded
domains, so they are rational with poles at roots of unity, with integral
coefficients. By Theorem~\ref{thm:K023}, $K_n\equiv B_n\pmod 3$. By
Theorem~\ref{thm:K024}, $B_n$ has a simple pole at each primitive $M$-th root
$\zeta$, $M=2^{n-1}$, with limit a unit in $\Z[\zeta,1/2]$, hence a unit at
every prime above $3$, and no pole at the primitive $M3^r$-th roots, $r\ge1$.
Theorem~\ref{thm:transfer} with $p=3$, $F=K_n$, $G=B_n$ applies.
\end{proof}

\begin{corollary}\label{cor:kroneckerPH3}
The existence clause of PH3 for Kronecker coefficients \cite[Hypothesis
3.3.4]{GCT6} fails: for the inputs $((1^n),(1^n),(1^n))$, with
$\rank(x)=n^2+5n$, every positive rational form of the generating function has
modular index at least $2^{n-1}$. The relaxed PH3 of \cite[Appendix,
p.~20]{BOR}, which bounds the modular index of the minimal positive form by a
polynomial in the largest height $h=n$, fails as well.
\end{corollary}

\begin{proof}
The minimal positive form is a positive form, and $2^{n-1}$ exceeds every
polynomial in $n$, hence in $n^2+5n$, for large $n$.
\end{proof}

\section{Plethysm coefficients: the modular index}\label{sec:plethysm}

Let $V=\C^n$ with coordinates $x_1,\dots,x_n$, $n\ge4$, $\nu=\gcd(n,3)$,
$d_0=n/\nu$, $\gamma=3/\nu$ (an odd integer) and $M=2^{n-1}$. Points of
$\Sym^3V^*$ are cubic forms, with $(a\cdot F)(x)=F(a^{-1}x)$. Let
$R=\C[\Sym^3V^*]^{SL_n}$ with its coefficient grading $R_D$.

\paragraph{The multiplicity algebra.}
If $R_D\ne0$, the scalar $\zeta I_n\in SL_n$, $\zeta^n=1$, acts on $R_D$ by
$\zeta^{\pm3D}$, so $n\mid3D$, that is $D=d_0s$. Writing $a=zh$ with
$h\in SL_n$ and $z^n=\det a$, the element $a\in GL_n$ acts on $R_{d_0s}$ by the
character $(\det a)^{\gamma s}$, and, as functions,
\begin{equation}\label{eq:detcharacter}
 P(a\cdot F)=(\det a)^{-\gamma s}P(F)\qquad(P\in R_{d_0s}).
\end{equation}
Choosing a volume vector $\varepsilon\in\Lambda^nV$, evaluation of an
intertwiner at $\varepsilon^{\otimes\gamma s}$ identifies the degree-$s$ part
of the multiplicity algebra $\bigoplus_s\Hom_{GL_n}((\det)^{\gamma s},
\Sym^{d_0s}(\Sym^3V))$ of \cite[Section 4.1.2, eq.~(4.5)]{GCT6} with $R_{d_0s}$,
compatibly with products. So the multiplicity algebra is
$T^{(n)}=\bigoplus_sR_{d_0s}$, the whole invariant ring with its degrees
divided by $d_0$, a finitely generated positively graded domain with
$T^{(n)}_0=\C$, and $H_n=\Hilb_{T^{(n)}}$. Let $X=\Spec T^{(n)}$. By
\eqref{eq:detcharacter}, $a\in GL_n$ acts on $X$ as the grading element
$(\det a)^{-\gamma}$; as $\det$ is onto $\C^*$ and $\gamma$ is odd, every
grading element arises.

\begin{proposition}\label{prop:chaincubic}
Let $F_n=x_1^2x_2+x_2^2x_3+\dots+x_{n-1}^2x_n+x_n^3$ and $p_n\in X$ its image.
\begin{enumerate}
\item $F_n$ defines a smooth hypersurface, and its $SL_n$-orbit is closed.
\item Every linear transformation preserving $F_n$ up to a scalar is a scalar
multiple of a diagonal matrix $\diag(a_1,a_1^{-2},a_1^4,\dots,a_1^{(-2)^{n-1}})$
with $a_1^{3M}=1$. These matrices form the stabilizer of $F_n$ in $GL_n$, a
cyclic group of order $3M$; the stabilizer in $SL_n$ is the scalar group
$\mu_\nu$, which acts trivially on cubic forms.
\item The stabilizer of $p_n$ in the grading group is exactly $\mu_M$.
\item $X$ is smooth at $p_n$, and for a primitive $M$-th root of unity $\zeta$
the fixed scheme $X^\zeta$ is smooth of dimension one at $p_n$.
\end{enumerate}
\end{proposition}

\begin{proof}
(1) The partial derivatives are $2x_1x_2$, $x_{i-1}^2+2x_ix_{i+1}$
($1<i<n$) and $x_{n-1}^2+3x_n^2$. At a common zero $x\ne0$, let $j$ be least
with $x_j\ne0$. If $j<n$, the $j$-th derivative gives $x_{j+1}=0$, and then
the $(j+1)$-st is $x_j^2\ne0$; if $j=n$, the last one is $3x_n^2\ne0$. So
$F_n$ is smooth. For the closed orbit, put $r_i=(1-(-1/2)^i)/3$ for $i<n$ and
$r_n=(1-r_{n-1})/3$; all are positive, and with $a_i=2e_i+e_{i+1}$ ($i<n$),
$a_n=3e_n$, one has $\sum_ir_ia_i=\one$. Fix an $SU_n$-invariant Hermitian
product on $\Sym^3V^*$; distinct monomials are orthogonal, having distinct
torus weights. The cubic $F=\sum_i(r_i/\|m_i\|^2)^{1/2}m_i$, where $m_i$ are the
chain monomials, is a zero of the moment map of $SU_n$: its diagonal part is the restriction of $\sum_ir_ia_i$ to the trace-free
diagonal matrices, which is zero because $\sum_ir_ia_i=\one$, and its off-diagonal part
vanishes because no two of the weights $a_i$ differ by a root $e_a-e_b$
(their differences have an entry $\pm2$, or the form $2e_i-e_{i+1}-e_{i+2}$, or
$2e_{n-1}-2e_n$). By the Kempf--Ness criterion \cite[Theorem 4, p.~1]{Losev}
the $SL_n$-orbit of $F$ is closed. The weights $a_i$ are linearly independent,
so a diagonal matrix maps $F$ to $F_n$, whose orbit is therefore closed.

(2) Half the Hessian matrix of $F_n$ is tridiagonal, with diagonal
$(x_2,x_3,\dots,x_n,3x_n)$ and off-diagonal entries $x_1,\dots,x_{n-1}$. Its
projective locus of rank at most two consists of
\begin{equation}\label{eq:hessianlocus}
 [e_1]\quad\text{and}\quad[ae_{n-1}+be_n]\ \text{with}\ a=0\ \text{or}\ a^2=3b^2 ,
\end{equation}
three collinear points and a fourth point off their line. Indeed, a tridiagonal
block whose off-diagonal entries are all nonzero has kernel of dimension at
most one, and the zero off-diagonal entries cut the matrix into such blocks; so
the rank is at least the number of nonzero $x_1,\dots,x_{n-1}$, and at most two
of them are nonzero. If exactly two, $x_i,x_j$ with $i<j$: when $j\ge i+2$ the
block on $\{i,i+1\}$ has determinant $-x_i^2$ (its first diagonal entry
$x_{i+1}$ vanishes) and the block of $x_j$ adds rank at least one; when $j=i+1$
and $i\ge2$, the singleton block $\{i-1\}$ has diagonal $x_i\ne0$, separate
from a block of rank at least two; when $i=1$, the block on $\{1,2,3\}$ has
determinant $-x_2^3\ne0$ for $n\ge5$, and for $n=4$ either this holds or
$x_4\ne0$ and the singleton $\{4\}$ adds one. If exactly one, $x_j$: for $2\le
j\le n-2$ the block $\{j,j+1\}$ has determinant $-x_j^2$ and the singleton
$\{j-1\}$ has diagonal $x_j$; for $j=1$ the block $\{1,2\}$ has rank two and
$x_n\ne0$ would add the last two diagonal entries, leaving $[e_1]$; for
$j=n-1$, with $a=x_{n-1}$, $b=x_n$, the singleton $\{n-2\}$ has entry $a$ and
the block $\left(\begin{smallmatrix}b&a\\a&3b\end{smallmatrix}\right)$ has rank
one exactly when $a^2=3b^2$. If none, the point is $[e_n]$, of rank two.

A linear transformation $T$ with $F_n\circ T=cF_n$ transforms the Hessian by
congruence, so it preserves \eqref{eq:hessianlocus}, hence fixes $[e_1]$, the
point off the line of the other three. In coordinate forms $\ell_i=(Tx)_i$
this says $\ell_1=ax_1+L(x_2,\dots,x_n)$ with $a\ne0$, while $\ell_2,\dots,
\ell_n$ do not involve $x_1$. The coefficient of $x_1^2$ in
$F_n(\ell)=cF_n(x)$ gives $a^2\ell_2=cx_2$; that of $x_1$ gives $2aL\ell_2=0$,
so $L=0$. Setting $x_1=0$, $(\ell_2,\dots,\ell_n)$ preserves up to a scalar the
chain cubic in $x_2,\dots,x_n$, with $\ell_2$ a multiple of $x_2$. If $n-1\ge4$
we repeat the argument; it remains to treat the ternary cubic
$G=y^2z+z^2w+w^3$ with $\ell=(by,uy+A,vy+B)$, $A,B$ independent linear forms
in $z,w$. The coefficient of $y$ in $G(\ell)$ is $vA^2+2uAB+3vB^2$; it
vanishes, so $u=v=0$. The coefficient of $y^2$ gives $b^2A=cz$; writing
$B=\alpha z+\beta w$ with $\beta\ne0$, the coefficient of $zw^2$ is
$3\alpha\beta^2=0$, so $\alpha=0$. Hence every such $T$ is diagonal. Dividing $T$ by a cube root of $c$, we may
take $c=1$. A diagonal $\diag(a_1,\dots,a_n)$ preserves $F_n$ exactly when $a_i^2a_{i+1}=1$
and $a_n^3=1$, that is $a_i=a_1^{(-2)^{i-1}}$ and $a_1^{3M}=1$; its
determinant is $a_1^S$ with $S=\sum_{i=0}^{n-1}(-2)^i=(1-(-2)^n)/3$, which is
odd and congruent to $n$ modulo three. So $\gcd(S,3M)=\nu$, and the kernel of
$\det$ has order $\nu$; when $\nu=3$ its elements are the scalars
$\omega I$, $\omega^3=1$, as $(-2)^i\equiv1\pmod3$, and these act trivially on
cubic forms.

(3) Since each fibre of the quotient map contains one closed orbit, a grading
element fixes $p_n$ exactly when it is $(\det a)^{-\gamma}$ for some $a$ in the
stabilizer of $F_n$. With $a_1$ of order $3M$, $(\det a)^{-\gamma}=
a_1^{-\gamma S}$ has order $3M/\gcd(3M,\gamma S)$. As $S$ is odd,
$\gcd(3M,\gamma S)=\gcd(3,\gamma S)$, which is $\gcd(3,3S)=3$ if $\nu=1$ and
$\gcd(3,S)=3$ if $\nu=3$ (then $S\equiv n\equiv0\pmod3$). So the grading
stabilizer of $p_n$ is $\mu_M$.

(4) The group $SL_n/\mu_\nu$ acts on $\Sym^3V^*$ and has trivial stabilizer at
$F_n$, whose orbit is closed. By Luna's slice theorem \cite[Section III.1,
Th\'eor\`eme du slice \'etale and Remarque 1$^\circ$, pp.~97--98]{Luna}, $X$ is smooth at $p_n$ with tangent space
$\Sym^3V^*/\mathfrak{sl}_n\cdot F_n$. Let $A=\diag(\zeta,\zeta^{-2},\dots,
\zeta^{(-2)^{n-1}})$ with $\zeta$ a primitive $M$-th root of unity; it fixes
$F_n$, and acts on $X$ as the grading element $(\det A)^{-\gamma}=\zeta^{-\gamma
S}$, a primitive $M$-th root of unity, whose fixed scheme is $X^\zeta$. The
monomials fixed by $A$ are exactly the $n$ chain monomials, by the valuation
argument at the end of the proof of Proposition~\ref{prop:graphsupport}; and
$\dim\mathfrak{sl}_n^A=n-1$, as $A$ has simple spectrum. Taking invariants is
exact, so $\dim(T_{p_n}X)^A=1$, and the fixed scheme is smooth of dimension one
at $p_n$ (proof of Proposition~\ref{prop:fixedcomponentpole}).
\end{proof}

\begin{theorem}\label{thm:plethysmPH3}
Let $n\ge4$ and $M=2^{n-1}$.
\begin{enumerate}
\item The support of the fixed scheme $X^{\mu_M}$ is the closure $Y$ of the
grading orbit of $p_n$, and $X$ is smooth along $Y\setminus\{0\}$.
\item $H_n$ has a simple pole at every primitive $M$-th root of unity $\zeta$,
with
\begin{equation}\label{eq:plethysmresidue}
 \mathcal C_n(\zeta):=\lim_{t\to\zeta}(1-t^M)H_n(t)
 =\frac1{\det\bigl(1-\zeta\mid N^*_{Y/X,\eta_Y}\bigr)}\ne0 .
\end{equation}
\item If $\eta$ is a root of unity whose order $q$ has two-part at least $M$,
then $H_n$ has a pole at $\eta$ if and only if $q=M$.
\item In every expression $H_n(t)=P(t)/\prod_j(1-t^{a_j})^{b_j}$ with
$P\in\C[t]$ and $a_j,b_j>0$, some $a_j$ is divisible by $M$.
\end{enumerate}
\end{theorem}

\begin{proof}
Let $p\ne0$ be a closed point fixed by a grading element of order $2^a\ge M$,
and $F$ a cubic in its closed orbit. Lemma~\ref{lem:lift} (with the odd
exponent $-\gamma$) gives a diagonal $A$ of two-power order fixing $F$ with
$\ord(\det A)=2^a$; Lemma~\ref{lem:span} gives the span condition; and
Theorem~\ref{thm:fullcubicsupport} gives $2^a\le M$, so $a=n-1$, and $F$ is
equivalent to $F_n$ under a permutation and a diagonal substitution. So
$p\in GL_n\cdot p_n$, which is the grading orbit of $p_n$. Conversely that
orbit is fixed by $\mu_M$ (Proposition~\ref{prop:chaincubic}(3)). The closure
$Y$ adds only the origin, so the support of $X^{\mu_M}$ is $Y$, a curve, and
$X$ is smooth along $Y\setminus\{0\}$ by Proposition~\ref{prop:chaincubic}(4)
and translation. Proposition~\ref{prop:fixedcomponentpole} with $\delta=1$ and
Lemma~\ref{lem:orbitclosure} ($e_Y=1$) give (2). For (3): if $q=2^ab$, $b$ odd,
$a\ge n-1$, and $q\ne M$, a nonzero point fixed by $\eta$ is fixed by its
two-primary part, of order $2^a\ge M$, hence lies on the grading orbit of
$p_n$, whose grading stabilizer is $\mu_M\not\ni\eta$. So $X^\eta$ is supported
at the origin, and Lemma~\ref{lem:isolated} applies. (4) follows from (2).
\end{proof}

\begin{corollary}\label{cor:plethysmPH3}
For $n\ge4$, every positive rational form of $H_n$ has modular index at least
$2^{n-1}$, while $\dim GL_m+\operatorname{ht}(\lambda)+\operatorname{ht}(\pi)
=m^2+n+1=O(n^6)$ and $\rank(x)=2n+m+1$. The existence clause of
\cite[Hypothesis 1.6.7]{GCT6}, and so also its specific clause on the minimal
positive form, fails for plethysm coefficients; so does
\cite[Hypothesis 3.3.4]{GCT6} for the plethysm problem.
\end{corollary}

\section{Plethysm coefficients: the positivity and saturation indices}\label{sec:indices}

We keep the notation of Section~\ref{sec:plethysm}, with $n\ge4$, and put
$\zeta=e^{2\pi i/M}$, $\mathcal C_n(\rho)=\lim_{t\to\rho}(1-t^M)H_n(t)$ and
$A_n(\rho)=\mathcal C_n(\rho)/M$ for primitive $M$-th roots $\rho$. Since
$H_n$ has rational coefficients, $\mathcal C_n(\zeta^a)$ is the Galois
conjugate of $\mathcal C_n(\zeta)$ under $\zeta\mapsto\zeta^a$, and
$A_n(\zeta^{-a})=\overline{A_n(\zeta^a)}$. Define the $M$-periodic
\emph{dyadic wave}
\begin{equation}\label{eq:wave}
 W_n(r)=\sum_{\substack{1\le a<M\\a\ \mathrm{odd}}}A_n(\zeta^a)\zeta^{-ar}
 =\frac1M\Tr_{\Q(\zeta)/\Q}\bigl(\mathcal C_n(\zeta)\zeta^{-r}\bigr)\in\Q .
\end{equation}

\subsection{The wave and the indices}

By \cite[Theorem 1.6.1(b), p.~21]{GCT6}, $f_n$ is a quasi-polynomial, that is,
$H_n$ is a proper rational function with poles at roots of unity. So
\begin{equation}\label{eq:fourier}
 f_n(s)=\sum_\rho P_\rho(s)\rho^{-s}\qquad(s\ge0),
\end{equation}
the sum over the poles $\rho$ of $H_n$, with $P_\rho$ a polynomial of degree
one less than the pole order. At a simple pole $\rho$ of order $M$,
$1-t^M=(1-t/\rho)\prod_{\rho'\ne\rho}(1-t/\rho')$ over the other $M$-th roots
of unity $\rho'$, and $\prod_{\rho'\ne\rho}(1-\rho/\rho')=M$; so near $\rho$,
$H_n(t)=(\mathcal C_n(\rho)/M)/(1-t/\rho)+O(1)$, and $P_\rho=A_n(\rho)$ is
constant.

\begin{proposition}\label{prop:wave}
Let $L$ be the least common multiple of the orders of the poles of $H_n$, a
period of $f_n$, and write $f_n(s)=q_r(s)$ for $s\equiv r\pmod L$. Then
$v_2(L)=n-1$, and for every residue $r$ and every $X$,
\begin{equation}\label{eq:waveidentity}
 q_r(X)-q_{r+L/2}(X)=2W_n(r).
\end{equation}
Consequently $f_n(x)\ge2W_n(x)$ for every integer $x\ge p(f_n)$ and every
integer $x\ge s(f_n)+1$. In particular, if $2W_n(x)>f_n(x)$ for an integer
$x\ge0$, then $p(f_n)>x$ and $s(f_n)\ge x$.
\end{proposition}

\begin{proof}
By Theorem~\ref{thm:plethysmPH3}(2),(3), every pole order has two-part at most
$M$, and $M$ occurs; so $v_2(L)=n-1$. By \eqref{eq:fourier},
$q_r(X)=\sum_\rho P_\rho(X)\rho^{-r}$. If the order of $\rho$ has two-part at
most $M/2$, it divides $L/2$ and $\rho^{-L/2}=1$; if it has two-part $M$, then
$\rho=\zeta^a$ with $a$ odd (Theorem~\ref{thm:plethysmPH3}(3)) and
$\rho^{-L/2}=-1$, as $L/M$ is odd. Subtracting the two expansions gives
\eqref{eq:waveidentity}. If $p=p(f_n)$, every constituent polynomial of
$n\mapsto f_n(n+p)$, namely $X\mapsto q_r(X+p)$, has nonnegative coefficients,
so $q_r(y)\ge0$ for all real $y\ge p$ and all $r$. If $s=s(f_n)$, every
constituent $X\mapsto q_r(X+s)$ is either identically zero or positive at every
integer $X\ge1$, so $q_r(y)\ge0$ for every integer $y\ge s+1$. At such $x$, take
$r\equiv x\pmod L$ in \eqref{eq:waveidentity}: $f_n(x)-2W_n(x)=
q_{r+L/2}(x)\ge0$.
\end{proof}

\subsection{An exact product for the residues}

Write $w_i=(-2)^i$ for $0\le i\le n-2$ and $w_{n-1}=0$, and form the multisets
\begin{equation}\label{eq:Sn}
\begin{gathered}
 \mathcal E_n=\{\{0,1,k\}:2\le k<n\}\cup\{\{0,2,3\}\},\\
 \mathcal S_n=\{w_i:0\le i\le n-2\}\ \sqcup\ \{w_i+w_j+w_k:0\le i<j<k<n,\
 \{i,j,k\}\notin\mathcal E_n\},
\end{gathered}
\end{equation}
with repetitions retained. Then $|\mathcal S_n|=(n-1)+\binom n3-(n-2)-1=
\binom n3$, and no element of $\mathcal S_n$ is divisible by $M$: a singleton
$w_i$ ($i\le n-2$) has valuation $i<n-1$, and in a sum of distinct $w$'s the
smallest valuation among the nonzero terms occurs once.

\begin{proposition}\label{prop:product}
For every odd $a$,
\begin{equation}\label{eq:product}
 \mathcal C_n(\zeta^a)=\prod_{b\in\mathcal S_n}\bigl(1-\zeta^{a\nu b}\bigr)^{-1}.
\end{equation}
\end{proposition}

\begin{proof}
By Galois conjugation it suffices to take $a=1$. Put $\eta=\zeta^\nu$ and
$\alpha_i=\eta^{w_i}$, and let $D=\diag(\alpha_0,\dots,\alpha_{n-1})$, which
fixes $F_n$ (its chain monomials have exponents $2w_i+w_{i+1}\equiv0$ and
$3w_{n-1}=0$). On coordinate functions of $\Sym^3V^*$, that is on
$\Sym(\Sym^3V)$, $D$ acts through the character of $\Sym^3V$, and on invariants
of coefficient degree $d_0s$ by $(\det D)^{\gamma s}$. Here
$(\det D)^\gamma=\zeta^{3\sum_iw_i}$ and $3\sum_iw_i=1-(-2)^{n-1}\equiv1
\pmod M$; so $D$ acts on $T^{(n)}$ as the grading element $\zeta$. By the
tangent space computation in Proposition~\ref{prop:chaincubic}(4), the
cotangent space of $X$ at $p_n$ has the character $\Sym^3V-(V\otimes V^*-1)$
as a representation of $\langle D\rangle$, and the conormal space of $Y$ is
obtained by removing the trivial character of the fixed direction:
\begin{equation}\label{eq:conormal}
 N^*=\Sym^3V-V\otimes V^* .
\end{equation}
With $h_j$, $e_j$ the complete and elementary symmetric functions in the
$\alpha_i$, the identity $\prod_i(1+t\alpha_i)=\prod_i(1-t^2\alpha_i^2)/
\prod_i(1-t\alpha_i)$ gives $h_3=e_3+h_1\sum_i\alpha_i^2$. The multisets
$\{\alpha_i^2\}$ and $\{\alpha_i^{-1}\}$ differ only in that $\alpha_0^{-1}$ is
replaced by $1$, because $\alpha_i^2=\alpha_{i+1}^{-1}$ for $i<n-1$ and
$\alpha_{n-1}=\alpha_{n-2}^2=1$. So \eqref{eq:conormal} is
$e_3+h_1-\alpha_0^{-1}h_1$. The trivial weight $\alpha_{n-1}$ in $h_1$ cancels
the term $i=0$ of $\alpha_0^{-1}h_1$. For $2\le k<n$ the subtracted weight
$w_k-w_0$ equals the weight of the triple $\{0,1,k\}$, as $w_0+w_1=-1$; the
last subtracted weight $w_1-w_0=-3$ equals that of $\{0,2,3\}$ ($1+4-8$ for
$n\ge5$; $1+4+0\equiv-3\pmod8$ for $n=4$). These triples are distinct, and what
remains is the multiset $\mathcal S_n$. Equation \eqref{eq:plethysmresidue}
gives \eqref{eq:product}.
\end{proof}

\subsection{A dominant pair}

Put $r_\nu=3\nu\in\{3,9\}$, let $1\le a_*<M$ with $r_\nu a_*\equiv1\pmod M$
(so $a_*$ is odd and $a_*\not\equiv-a_*$), and
\begin{equation}\label{eq:ZE}
 Z_n=A_n(\zeta^{a_*}),\qquad
 E_n=\sum_{\substack{a\bmod M\ \mathrm{odd}\\a\not\equiv\pm a_*}}|A_n(\zeta^a)| .
\end{equation}

\begin{proposition}\label{prop:amplitude}
For every $n\ge4$,
\begin{equation}\label{eq:amplitude}
 \log_2|Z_n|\ge\binom n4-3\binom{n-1}3-(n-1)(n-2)-\binom n3-(n-1)
 \ge\frac{n^4}{24}-2n^3 .
\end{equation}
If $n\ge6$ and $E_n\le|Z_n|/4$, then for every integer $X$ there is an integer
$x$ with $X\le x\le X+8$ and $W_n(x)\ge|Z_n|/4$.
\end{proposition}

\begin{proof}
By \eqref{eq:product}, $Z_n=M^{-1}\prod_{b\in\mathcal S_n}(1-\theta^b)^{-1}$ with
$\theta=\zeta^{\nu a_*}$, so that $\theta^3=\zeta$. Every factor has absolute
value at most two. Consider the triples $i<j<k\le n-2$ of $\mathcal S_n$, those
not using the zero weight. Their weight $w=w_i+w_j+w_k$ is divisible by three
(each $w_i\equiv1\pmod3$) and $|w|<2^{k+1}$, so
$|1-\theta^w|=|1-\zeta^{w/3}|\le2\pi|w|/(3M)\le8\cdot2^k/M$, and
$-\log_2|1-\theta^w|\ge(n-1-k)-3$. Summing $n-1-k$ over all triples
$i<j<k\le n-2$ counts the choices of a fourth index in $\{k+1,\dots,n-1\}$ and
gives $\binom n4$; at most $n-2$ of these triples belong to $\mathcal E_n$, each
with a term at most $n-1$. The remaining at most $\binom n3$ factors contribute
at least $-1$ each, and $M^{-1}$ contributes $-(n-1)$. This proves the first
inequality of \eqref{eq:amplitude}. For the second use $\binom n4\ge
n^4/24-n^3/4$, $3\binom{n-1}3\le n^3/2$, $\binom n3\le n^3/6$ and
$(n-1)(n-2)+(n-1)\le n^2\le n^3/4$: the total cubic loss is at most
$\tfrac76n^3$.

Write $r_\nu a_*=1+kM$; then $1\le k<r_\nu$ and $\gcd(k,r_\nu)=1$, since
$kM\equiv-1\pmod{r_\nu}$. Thus $a_*/M=k/r_\nu+1/(r_\nu M)$. For fixed $X$, the
$r_\nu$ numbers $Z_n\zeta^{-a_*X}e^{-2\pi ikj/r_\nu}$, $0\le j<r_\nu$, are
equally spaced on the circle of radius $|Z_n|$, so one has real part at least
$\cos(\pi/r_\nu)|Z_n|\ge|Z_n|/2$. Replacing the phase $e^{-2\pi ikj/r_\nu}$ by
$\zeta^{-a_*j}$ changes it by a factor $e^{-2\pi ij/(r_\nu M)}$ and the real
part by at most $(2\pi/M)|Z_n|\le|Z_n|/4$ for $n\ge6$. So
$2\operatorname{Re}(Z_n\zeta^{-a_*(X+j)})\ge|Z_n|/2$ for this $j\le8$. The terms
$a=\pm a_*$ of \eqref{eq:wave} are this conjugate pair, and the others have
total absolute value at most $E_n\le|Z_n|/4$.
\end{proof}

\begin{corollary}\label{cor:conditional}
If $E_n\le|Z_n|/4$ for all sufficiently large $n$, then for every real
$0<c<1/4$ and all sufficiently large $n$, $p(f_n)>2^{cn}$ and $s(f_n)\ge2^{cn}$.
\end{corollary}

\begin{proof}
Take $X=\lceil2^{cn}\rceil$ and $x\le X+8$ as in
Proposition~\ref{prop:amplitude}. The value $f_n(x)$ is at most the dimension
of $\Sym^{d_0x}(\Sym^3V)$, so
$\log_2\max\{1,f_n(x)\}\le(m-1)\log_2(m+nx)=\frac c6n^4+O_c(n^3\log n)$, which
is eventually smaller than $\log_2(2W_n(x))\ge n^4/24-2n^3-1$, as $c/6<1/24$.
Proposition~\ref{prop:wave} gives $p(f_n)>x\ge2^{cn}$ and $s(f_n)\ge x$.
\end{proof}

\subsection{Reduction to an increment}

For $n\ge4$ write $M_n=2^{n-1}$, $b_n=3^{-1}\bmod M_n$, and for odd $a$
\begin{equation}\label{eq:Pn}
 P_n(a)=\prod_{b\in\mathcal S_n}|1-e^{2\pi iab/M_n}|^{-1},\qquad
 Q_n=\sum_{\substack{a\bmod M_n\ \mathrm{odd}\\a\ne\pm b_n}}\frac{P_n(a)}{P_n(b_n)} .
\end{equation}
By \eqref{eq:product}, $|A_n(\zeta^a)|=P_n(\nu a)/M_n$; multiplication by the
odd number $\nu$ permutes the odd residues, and $\nu a_*\equiv b_n$. So
$Q_n=E_n/|Z_n|$ and $P_n(b_n)=M_n|Z_n|$. For $n\ge5$ form the multiset of
$\binom{n-1}2$ residues
\begin{equation}\label{eq:Tn}
 \mathcal T_n=\{1\}\sqcup\{1+w_j+w_k:2\le j<k<n,\ (j,k)\ne(2,3)\}
 \sqcup\{2+w_k:3\le k<n\}\sqcup\{6\},
\end{equation}
and $J_n(a)=\prod_{b\in\mathcal T_n}|1-e^{2\pi iab/M_n}|^{-1}$.

\begin{proposition}\label{prop:recursion}
For $n\ge5$ and odd $a$, $P_n(a)=J_n(a)\,P_{n-1}(-a\bmod M_{n-1})$. The selected
sibling ratio satisfies
\begin{equation}\label{eq:sibling}
 \delta_n:=\frac{P_n(b_n+M_{n-1})}{P_n(b_n)}\le2^{-n^3/6+4n^2-107n/6+19},
 \qquad\delta_n\le2^{-n^3/12}\ (n\ge48).
\end{equation}
If for some $0<\rho<1$ and all $n\ge n_0$ and all odd $a$ with
$a\bmod M_{n-1}\notin\{\pm b_{n-1}\}$
\begin{equation}\label{eq:contraction}
 J_n(a)+J_n(a+M_{n-1})\le\rho J_n(b_n),
\end{equation}
then $Q_n\le\rho Q_{n-1}+2\delta_n$ for $n\ge n_0$, and $Q_n\to0$.
\end{proposition}

\begin{proof}
The weights $w_1,\dots,w_{n-1}$ are $-2$ times the weights of the chain in
$n-1$ variables, in the same order. All triples with smallest index at least one
lie in $\mathcal S_n$; the scaled $\mathcal S_{n-1}$ omits among them exactly
$\{1,2,k\}$ ($3\le k<n$) and $\{1,3,4\}$; and $\mathcal S_n$ has the extra
singleton $w_0=1$ and the triples $\{0,j,k\}$, $2\le j<k<n$, other than
$\{0,2,3\}$. Their weights are those of \eqref{eq:Tn}: $w_1+w_2+w_k=2+w_k$, and
$\{1,3,4\}$ has weight $-2-8+16=6$ ($n\ge6$) or $-2-8+0\equiv6\pmod{16}$ ($n=5$).
The count is $1+(\binom{n-2}2-1)+(n-3)+1=\binom{n-1}2$. Substituting
$e^{2\pi i a(-2b')/M_n}=e^{2\pi i(-a)b'/M_{n-1}}$ gives the recursion.

The residues $b_n$ and $b_n+M_n/2$ have the same smaller factor, and every even
weight gives the same factor at both; so $\delta_n=\prod_{b\in\mathcal O_n}
|\tan(\pi b_nb/M_n)|$ over the odd weights $\mathcal O_n=\{1\}\sqcup\{1+w_j+w_k\}$.
The bulk ones, $2\le j<k\le n-2$, $(j,k)\ne(2,3)$, number
$B=\binom{n-3}2-1$, are divisible by three with $|b|<2^{k+1}\le M_n$; as
$3b_n\equiv1$, $|\tan(\pi b_nb/M_n)|=|\tan(\pi b/(3M_n))|\le2\pi|b|/(3M_n)\le
8\cdot2^k/M_n$, using $|\sin u|\le|u|$ and $\cos u\ge1/2$ on $[-\pi/3,\pi/3]$.
The sum of $n-1-k$ over them is $\binom{n-2}3-(n-4)$. The other $n-2$ odd
factors are at most $M_n/2$, an odd residue being at distance at least one from
$M_n/2$, so the cosine is at least $\sin(\pi/M_n)\ge2/M_n$. Hence
$\log_2\delta_n\le3B-\binom{n-2}3+(n-4)+(n-2)^2=-n^3/6+4n^2-107n/6+19$, which
is at most $-n^3/12$ for $n\ge48$.

Group the odd residues modulo $M_n$ by their reduction modulo $M_{n-1}$. The
two selected parents $\pm b_{n-1}$ have the selected children $\pm b_n$ and the
children $\pm b_n+M_{n-1}$, of total normalized mass $2\delta_n$ (conjugation).
Every other parent has two children whose normalized mass is at most $\rho$
times its own in $Q_{n-1}$, by the recursion and \eqref{eq:contraction}, since
$P_n(b_n)=J_n(b_n)P_{n-1}(b_{n-1})$. So $Q_n\le\rho Q_{n-1}+2\delta_n$, and as
$\delta_n\to0$, $Q_n\to0$.
\end{proof}

\subsection{The increment estimates}

\begin{lemma}\label{lem:sparse}
Let $1\le r\le3$, integers $1\le q_1<\dots<q_r$, signs $\epsilon_i$,
$\kappa\in\Z$, $J\ge1$ odd, $R=\lceil\log_2(16J)\rceil$, and
$x=\kappa/3+(J/3)\sum_i\epsilon_i2^{-q_i}$. Then $\|x\|_{\mathbb R/\Z}\ge
1/(6\cdot2^{q_1+(r-1)R})$. In particular, if $q_1\le R$, then
$2^{-3R-1}\le|1-e^{2\pi ix}|\le2$.
\end{lemma}

\begin{proof}
Keep the initial block of terms ending before the first gap $q_{i+1}-q_i>R$
(all terms if there is none), ending at $q_t$. Its numerator over $3\cdot
2^{q_t}$ is odd (the last term is odd, the earlier ones and $\kappa2^{q_t}$
even), so it is at distance at least $1/(3\cdot2^{q_t})$ from $\Z$. The tail
is less than $(2J/3)2^{-q_{t+1}}\le J2^{-R}/(3\cdot2^{q_t})\le
1/(48\cdot2^{q_t})$. So $\|x\|\ge1/(6\cdot2^{q_t})$, and $q_t\le q_1+(r-1)R$.
Finally $|1-e^{2\pi ix}|\ge4\|x\|$.
\end{proof}

Let $l=n-1$, $M=2^l$, and for odd $J$ let $b_J\equiv J\cdot3^{-1}\pmod M$; then
$\eta_J=\zeta^{b_J}$ satisfies $\eta_J^3=\zeta^J$ and
\begin{equation}\label{eq:cuberoot}
 \eta_J=\omega_Je^{2\pi iJ/(3M)},\qquad\omega_J^3=1,\qquad
 \omega_J=1\iff3\mid J\quad(J<M),
\end{equation}
since $3b_J=J+k_JM$ with $k_J\in\{0,1,2\}$, $\omega_J=e^{2\pi ik_J/3}$, and
$k_J=0$ only if $3\mid J$. For $|y|\le1/24$,
$|\log_2(\sin(\pi y)/(\pi y))|\le8y^2$, by $1-u^2/6\le\sin u/u\le1$.

\begin{proposition}[near increments]\label{prop:near}
For odd $1\le j<M$, with $a_j\equiv j\cdot3^{-1}\pmod M$ and
$R=\lceil\log_2(16j)\rceil$, and $n\ge5$,
\begin{equation}\label{eq:near}
 \Bigl|\log_2\frac{J_n(a_j)}{J_n(b_n)}+\binom{n-1}2\log_2j
 -\mathbf1_{3\mid j}\binom n2\Bigr|\le100nR^2 .
\end{equation}
For $\delta=10^{-5}$ and odd $3\le j\le2^{\delta n}$ this gives
$\log_2(J_n(a_j)/J_n(b_n))\le-n^2\log_2j/12$.
\end{proposition}

\begin{proof}
The weights of $\mathcal T_n$ are: the \emph{bulk} weights $1+w_u+w_v$
($2\le u<v\le l-1$, $(u,v)\ne(2,3)$), the weights $2+w_k$ ($3\le k\le l-1$) and
$6$, all divisible by three, $B_3=(n-1)(n-4)/2$ in number; and the
\emph{exceptional} weights $1$, $2$ and $1+w_k$ ($2\le k\le l-1$), $n-1$ in
number. Give each weight its largest active dyadic index $k$ and put $q=l-k$
(the weights $1$ and $2$ get $q=l$ and $q=l-1$). The exceptional weights have
each $q=1,\dots,l$ exactly once, so $B_3+(n-1)=\binom{n-1}2$ and the exceptional
$q$ sum to $\binom n2$. A weight divisible by three has absolute value between
$2^{k-2}$ and $2^{k+1}$ (for $1+w_u+w_v$ use $2^v-2^{v-1}-1\ge2^{v-2}$, $v\ge3$);
for each $q$ there are at most $2n$ of them. An exceptional weight has absolute
value between $2^{k-1}$ and $3\cdot2^{k-1}$. By \eqref{eq:cuberoot}, the target
root is $\omega_je^{2\pi ij/(3M)}$. A factor is \emph{low} if $q>R$ and
\emph{high} otherwise. Lemma~\ref{lem:sparse} applies to every factor (each
weight is a signed sum of at most three powers of two): a high target factor is
at least $2^{-3R-1}$, and at least $2^{-2R-1}$ for an exceptional one; all
factors are at most two.

For a low factor divisible by three there is no cube-root shift: its target
phase is $j$ times its reference phase and is at most $(2j/3)2^{-q}<1/24$ in
absolute value; by the sine estimate the logarithm of the factor ratio differs
from $\log_2j$ by at most $16j^24^{-q}$; summing over at most $2n$ weights for
each $q>R$ gives less than $n$, as $j^24^{-R}\le1/256$. For a high factor
divisible by three, the reference phase lies between $2^{-q}/12$ and
$2\cdot2^{-q}/3<1/2$ in absolute value, so the reference factor is at least
$2^{-q}/3$, and the logarithm of the ratio is at most $5R$ in absolute value;
there are at most $2nR$ of them, and restoring their $\log_2j$ main terms costs
at most $2nR^2$. So the factors divisible by three contribute $-B_3\log_2j$ to the
logarithm of the reciprocal product, with error at most $20nR^2$.

Every exceptional reference factor lies between $1/3$ and two: its cube-root
shift is nontrivial and its phase perturbation is at most $1/4$. If $3\nmid j$,
every low exceptional target factor has the same kind of shift, perturbed by
less than $1/24$, so its logarithmic ratio is at most three; the at most $R$
high ones are at most $3R$ each; the total is at most $10nR^2$. If $3\mid j$,
the target shift vanishes; a low exceptional target factor lies between
$(2/3)j2^{-q}$ and $4j2^{-q}$, so its reciprocal ratio has logarithm $q-\log_2j$
up to four. Restoring the $q$ of the at most $R$ high factors costs at most
$R^2$, their actual ratios at most $3R^2$, discarding the $-\log_2j$ terms at
most $nR$, and the low errors at most $4n$; the exceptional contribution is
$\binom n2$ up to $10nR^2$. Altogether
$\log_2(J_n(a_j)/J_n(b_n))=-B_3\log_2j+\mathbf1_{3\mid j}\binom n2+\mathcal E$ with
$|\mathcal E|\le30nR^2$, and replacing $B_3$ by $\binom{n-1}2$ costs
$(n-1)\log_2j\le nR$.

For the consequence, put $t=\log_2j>3/2$, so $R\le t+5$. With
$\binom{n-1}2\ge n^2/2-2n$ and $\binom n2\le n^2/2$, \eqref{eq:near} gives at
most $-(t-1)n^2/2+n(2t+100(t+5)^2)\le-tn^2/6+2502nt^2$, using $t-1\ge t/3$ and
$t+5\le5t$; and $t\le n/10^5$ with $2502/10^5<1/12$.
\end{proof}

\begin{proposition}[far increments]\label{prop:far}
Let $n\ge5$, $l=n-1$, $M=2^l$, $a$ odd, $j\equiv3a\pmod M$ with $1\le j<M$, and
$J=\min\{j,M-j\}$. Then
\begin{equation}\label{eq:far}
 \log_2\frac{J_n(a)}{J_n(b_n)}\le7n^2-\frac{(\log_2J-8)_+^3}{24}.
\end{equation}
In particular, for $\delta>0$, $J\ge2^{\delta n}$ and
$n\ge\max\{5,16/\delta,2688/\delta^3\}$, the ratio is at most
$2^{-\delta^3n^3/384}$.
\end{proposition}

\begin{proof}
Consider the \emph{active pair} factors $1-\exp(2\pi ia(1+w_j+w_k)/M)$,
$2\le j<k<l$; the increment contains all but $(j,k)=(2,3)$. Every phase is
nonzero modulo one, its numerator being odd. Put $\theta=a/M$,
$y_{jk}=\theta(1+w_j+w_k)$ modulo one, and for $0\le u\le l-1$ let $r=l-u$ and
$N_u$ be the number of active pairs with $\|y_{jk}\|<2^{-u}$.

The map sending a subset of $\{0,\dots,l-1\}$ to the sum of its $w_i$ modulo
$2^l$ is injective (a difference of two subset sums has a unique least index
with coefficient $\pm1$, of valuation less than $l$). Call an index early if
$j<r$ and late otherwise; there are $e=(r-2)_+$ early indices. There are at
most $\binom e2$ early--early pairs. For a fixed early $j$, the mixed numerators
$a(1+w_j+w_k)$ are congruent modulo $2^r$ and distinct modulo $M$; the condition
$\|y_{jk}\|<2^{-u}$ puts their centred representatives in $(-2^r,2^r)$, which
holds at most three of one class; so there are at most $3e$ good mixed pairs.
The late--late numerators are all congruent to $a$ modulo $2^r$ and distinct,
so at most three are good. Hence
\begin{equation}\label{eq:Nu}
 N_u\le\binom r2+3r+3 .
\end{equation}

Let $\tau=\|3\theta\|=J/M$ and suppose $\tau>3\cdot2^{-u}$. Partition the early
indices into $h=\lfloor e/2\rfloor$ consecutive pairs $(2,3),(4,5),\dots$. For
two of them, $(j,j+1)$ and $(k,k+1)$, the identity
$y_{jk}+y_{j+1,k}+y_{j,k+1}=3\theta\pmod1$ (from $w_{j+1}=-2w_j$) shows that
the three phases cannot all be within $2^{-u}$ of an integer; so one of the four
pairs of the block is bad. The $\binom h2$ blocks are disjoint and do not
contain $(2,3)$. So under this condition
$N_u\le\binom r2+3r+3-\binom{\lfloor(r-2)_+/2\rfloor}2$.

For a nonzero phase $y$ with denominator dividing $M$, $\|y\|\ge1/M$ and
$\log_2|1-e^{2\pi iy}|^{-1}\le-\log_2\|y\|\le\sum_{u=0}^{l-1}
\mathbf1_{\{\|y\|<2^{-u}\}}$. The other $2n-4$ increment factors have
logarithmic reciprocals at most $l$ each, less than $2n^2$ in total. Let
$K=(\lfloor\log_2J\rfloor-2)_+$; for $1\le r\le K$, $J>3\cdot2^r$, so the
condition holds at $u=l-r$. Summing,
\[
 \log_2J_n(a)\le\binom n3+\tfrac32l(l+1)+3l+2n^2
 -\sum_{r=1}^K\binom{\lfloor(r-2)_+/2\rfloor}2 ,
\]
where the first four terms are at most $n^3/6+4n^2$ for $n\ge5$, and the sum is
at least $(K-5)_+^3/24\ge(\log_2J-8)_+^3/24$, since for $r\ge6$ the summand is
at least $(r-3)(r-5)/8\ge(r-5)^2/8$.

For the selected increment, the $B=\binom{n-3}2-1$ bulk active pairs
(divisible by three) have factor at most $8\cdot2^k/M$ at $a=b_n$, as in the
proof of Proposition~\ref{prop:recursion}, and all other factors are at most
two. So
\[
 \log_2J_n(b_n)\ge\binom{n-2}3-(n-4)-2B-\binom{n-1}2
 =\frac{n^3}6-3n^2+\frac{71n}6-11\ge\frac{n^3}6-3n^2 .
\]
Subtracting gives \eqref{eq:far}. If $\log_2J\ge\delta n\ge16$, then
$(\log_2J-8)^3/24\ge\delta^3n^3/192$, and $7n^2\le\delta^3n^3/384$ for
$n\ge2688/\delta^3$.
\end{proof}

\begin{theorem}\label{thm:indices}
$E_n/|Z_n|\to0$. Consequently, for every real $0<c<1/4$ and all sufficiently
large $n$,
\[
 p(f_n)>2^{cn},\qquad s(f_n)\ge2^{cn},
\]
and $p(f_n)>2^{n/8}$, $s(f_n)\ge2^{n/8}$ for every $n\ge2\cdot10^{80}$.
\end{theorem}

\begin{proof}
For an odd residue $a\ne\pm b_n$, let $j$ be the least positive residue of $3a$
up to sign, $1<j\le M/2$ (and $j=1$ exactly for $a=\pm b_n$); as
$J_n(a)=J_n(-a)$, Proposition~\ref{prop:near} covers $j\le2^{\delta n}$ with
ratio at most $2^{-n^2\log_23/12}$, and Proposition~\ref{prop:far} the rest with
ratio at most $2^{-\delta^3n^3/384}$, for $n$ large. So the two children of any parent other than $\pm b_{n-1}$ have normalized increments of
sum at most $1/2$ eventually: \eqref{eq:contraction} holds with $\rho=1/2$, and
Proposition~\ref{prop:recursion} gives $Q_n=E_n/|Z_n|\to0$.
Corollary~\ref{cor:conditional} gives the first assertion.

For the explicit bound, let $N=10^{20}$. For $n\ge N$ both ratios are at most
$2^{-n^2/8}$ ($\log_23>3/2$; and $\delta^3n^3/384\ge n^2/8$ for
$n\ge48/\delta^3$), so \eqref{eq:contraction} holds with $\rho=1/2$ for all
$n\ge N$. Every factor of $P_n(a)$ is at least $4/M_n$, so
$P_n(a)\le2^{(n-3)\binom n3}$, while $P_n(b_n)=M_n|Z_n|\ge2^{-2n^3}$ by
\eqref{eq:amplitude}; hence $Q_n\le2^{n-1+(n-3)\binom n3+2n^3}\le2^{n^4}$.
Iterating from $N$, with $\delta_k\le2^{-k^3/12}$,
\[
 Q_n\le2^{-(n-N+1)}Q_{N-1}+2\sum_{k=N}^n2^{-(n-k)}\delta_k
 \le2^{-n+N-1+N^4}+4\cdot2^{-N^3/12}\le\frac14\qquad(n\ge2N^4).
\]
For such $n$ take $X=\lceil2^{n/8}\rceil$ in Proposition~\ref{prop:amplitude}:
some $x\in[X,X+8]$ has $2W_n(x)\ge2^{n^4/24-2n^3-1}$. As $\log_2n\le n/16$ for
$n\ge128$, $2^{n/8}\ge n^2$; also $m\le n^3/5$ for $n\ge16$ and $x\le2^{n/8+1}$,
so $m+nx\le4n2^{n/8}$ and
\[
 \log_2\max\{1,f_n(x)\}\le(m-1)\log_2(m+nx)\le\frac{n^3}5\Bigl(\frac n8+
 \log_2n+2\Bigr)\le\frac{3n^4}{80}+\frac{2n^3}5<\frac{n^4}{24}-2n^3-1
\]
for $n\ge1024$. Proposition~\ref{prop:wave} gives $p(f_n)>x\ge2^{n/8}$ and
$s(f_n)\ge x$.
\end{proof}

\begin{corollary}\label{cor:PH2SH}
Hypotheses 1.6.5(a) (SH) and 1.6.6 (PH2) of \cite{GCT6} fail for the plethysm
coefficients of the ray $\mu=(3)$, $\lambda=(d_0)$, $\pi=((3/\nu)^n)$: the
indices exceed every polynomial in $h=m^2+n+1=\Theta(n^6)$. Hypotheses
3.3.2(a) and 3.3.3 fail for the plethysm problem with $\rank(x)=2n+m+1$.
\end{corollary}

\section{Restriction from a torus}\label{sec:torus}

The subgroup restriction problem of \cite[Section 3.4]{GCT6} allows $H$ to be a
torus and $\rho:H\to G$ any explicit polynomial homomorphism
\cite[Section 3.4.1]{GCT6}; the rank of the input is the sum of the ranks of
$H$ and $G$ and the heights of the specified weights, and it does not see the
exponents of $\rho$. For $H=\Gm$ the multiplicity of the character $t\mapsto
t^{k}$ in the restriction of $\Sym^n\C^N$ is the number of monomials of degree
$n$ whose weight is $k$.

\begin{proposition}\label{prop:torusmodular}
For $N\ge2$ let $\rho_N:\Gm\to GL_2$, $\rho_N(t)=\diag(t^{N+1},t)$, let $\pi=(1)$
(the representation $\C^2$, stretched to $\Sym^n\C^2$) and $\lambda$ the
character $t\mapsto t^2$ (stretched to $t^{2n}$). Then
\[
 f_N(n)=\dim\Hom_{\Gm}\bigl(t^{2n},\Sym^n(\C^2)|_{\rho_N}\bigr)=\mathbf1_{N\mid n},
 \qquad \sum_{n\ge0}f_N(n)T^n=\frac1{1-T^N},
\]
and every expression of this series as $P(T)/\prod_i(1-T^{a_i})^{b_i}$ has some
$a_i$ divisible by $N$, while $\rank(x)=1+2+1+1=5$. Hypothesis 3.3.4 of
\cite{GCT6} fails for subgroup restriction as stated.
\end{proposition}

\begin{proof}
The monomial $y_1^ay_2^{n-a}$ has weight $(N+1)a+(n-a)=n+Na$, which equals
$2n$ exactly when $Na=n$; there is one such $a\in[0,n]$ if $N\mid n$ and none
otherwise. The series has a pole at every primitive $N$-th root of unity, and a
denominator must vanish there.
\end{proof}

Here $p(f_N)=s(f_N)=0$, so this example does not touch PH2 or SH.

\begin{proposition}\label{prop:torusindices}
For odd $N\ge3$ let $\sigma_N:\Gm\to GL_3$, $\sigma_N(t)=\diag(t,t^3,t^{N+1})$,
with $\pi=(1)$ and $\lambda=t^2$ as before. Then
\[
 g_N(n)=\#\{(b,c)\in\Z_{\ge0}^2:2b+Nc=n\},\qquad
 \sum_{n\ge0}g_N(n)T^n=\frac1{(1-T^2)(1-T^N)},
\]
and $p(g_N)=s(g_N)=N-2$, while $\rank(x)=1+3+1+1=6$. Hypotheses 3.3.2(a) and
3.3.3 of \cite{GCT6} fail for subgroup restriction as stated.
\end{proposition}

\begin{proof}
A monomial $y_1^ay_2^by_3^c$ of degree $n$ has weight $a+3b+(N+1)c=n+2b+Nc$,
which equals $2n$ exactly when $2b+Nc=n$, and then $a=n-b-c\ge0$ automatically.
This gives $g_N$ and its series. As $N$ is odd, $n-Nc$ is even exactly when
$c\equiv n\pmod2$, so $g_N(n)$ is the number of $c\equiv n\pmod2$ with
$0\le c\le n/N$. The period is $2N$. For $n=2Nq+r$, $0\le r<2N$, one finds
$g_N(n)=q+1$ if $r$ is even, and $g_N(n)=q+\mathbf1_{r\ge N}$ if $r$ is odd. So
the constituents are $g_r(X)=(X-r)/(2N)+\varepsilon_r$, with $\varepsilon_r=1$
for even $r$ and $\varepsilon_r=\mathbf1_{r\ge N}$ for odd $r$.

The constituents of $n\mapsto g_N(n+p)$ are $X\mapsto g_r(X+p)$, $0\le r<2N$,
with constant terms $(p-r)/(2N)+\varepsilon_r$ and slope $1/(2N)$. For odd
$r<N$ the constant term is $(p-r)/(2N)$, nonnegative only if $p\ge r$; the
largest such $r$ is $N-2$. For $p=N-2$ the other constant terms are positive:
$(N-2-r)/(2N)+1\ge1/2$ for even $r\le2N-2$, and $\ge(N-1)/(2N)$ for odd
$r\ge N$. So $p(g_N)=N-2$. A constituent $X\mapsto g_r(X+s)$ is not
identically zero and is increasing, so it is positive at every integer $X\ge1$
exactly when $(1+s-r)/(2N)+\varepsilon_r>0$; for odd $r<N$ this means $s\ge r$,
and for $s=N-3$ it fails at $r=N-2$; for $s=N-2$ all hold. So $s(g_N)=N-2$.
\end{proof}

Neither example touches the density clause 3.3.2(b). In both, the index is of
the size of the exponent $N$ of the embedding, while the specification of a
torus character has bit length $O(\log N)$ \cite[Section 3.4.2, p.~56]{GCT6}.
So the bounds fail in $\rank(x)$ and also in the bit length of the input. For
the general subgroup restriction problem a bound would have to depend on the
exponents of $\rho$ themselves; with that dependence these two examples are
harmless. The Kronecker and plethysm families of
Sections~\ref{sec:kronecker}--\ref{sec:indices} have no such parameter: their
embeddings are fixed by $n$.

\section{Earlier work and the scope of the results}\label{sec:literature}

\paragraph{What we searched.}
On 25 November 2026 we ran three web searches: for counterexamples to
Mulmuley's hypotheses PH2, PH3 and SH on the stretching quasi-polynomials of
Kronecker and plethysm coefficients; for the poles and periods of stretched
rectangular Kronecker coefficients; and for the work of Kahle and
Micha\l{}ek on the quasi-polynomials of plethysm. They led to \cite{BOR} and
\cite{KM}. We also searched the Hexagon repository
(\url{https://hexagonmath.org}) for ``Mulmuley'', ``plethysm'',
``Kronecker'', ``geometric complexity'', ``positivity hypothesis'',
``saturation hypothesis'' and ``modular index''; these searches returned no
work on the hypotheses (the matches for ``plethysm'' and ``Kronecker'' concern
other questions). We then took the works that cite \cite{GCT6} (64) or
\cite{BOR} (52) in the OpenAlex index, 94 works in all, and searched the full
text of the 73 of them that are on arXiv for the names of the hypotheses and
for ``positivity index'', ``saturation index'', ``modular index'' and
``positive form''; the 21 others we judged by title only. Besides Mulmuley's
own papers and \cite{BOR}, the matches discuss the strict saturation property
of particular families of (reduced) Kronecker coefficients. Finally we read the
four texts discussed below. We did not search zbMATH or MathSciNet. This is a
bounded search, and we claim no priority.

\paragraph{Briand, Orellana and Rosas.}
In \cite{BOR}, Briand, Orellana and Rosas showed that the strict saturation
hypothesis of the first versions of GCT~VI (``strong SH'') fails for Kronecker
coefficients: for $\lambda=(k,k)$, $\mu=(k+1,k-1)$, $\nu=(2k-2i-2j,2i,2j)$
with $i>j>0$ and $k>2i+j$, the stretched coefficient is $N/2+1$ for even $N$ and $(N-1)/2$ for odd $N$, so it vanishes at $N=1$ and is positive for every $N>1$ \cite[Theorem 2.4, p.~7]{BOR}; they also gave
counterexamples to strong PH2 \cite[Section 2.4, p.~8]{BOR}. In their examples
the positivity index is at most one \cite[Section 2.5, p.~8]{BOR}. Mulmuley's
appendix \cite[pp.~18--20]{BOR} then relaxed SH, PH2 and PH3 to polynomial
bounds in the largest height. Our Theorem~\ref{thm:kronecker} refutes the
relaxed PH3 of that appendix; the relaxed SH(a) and PH2 for Kronecker
coefficients remain open as far as this paper is concerned, while their
analogues for plethysm fail by Theorem~\ref{thm:indices}. GCT~VI v4 (22 January
2009), the newest version on arXiv, states the relaxed hypotheses in the form
quoted in Section~\ref{sec:defs}.

\paragraph{Hardness, rectangular coefficients, and a report.}
Ikenmeyer, Mulmuley and Walter proved that deciding the positivity of Kronecker
coefficients is NP-hard \cite[Theorem 1.3, p.~7]{IMW}, and constructed
superpolynomially many triples whose Kronecker coefficient vanishes while some
stretched coefficient does not \cite[Theorem 1.7, p.~11]{IMW}; these are
failures of the strict saturation property, and they bear on the exact decision
problem, while the relaxed hypotheses were formulated for a relaxed decision
problem \cite[Theorem 3.3.6, p.~52]{GCT6}. Fischer and Ikenmeyer proved the
same NP-hardness, and $\#P$-hardness of computation, for plethysm coefficients
with any fixed inner degree $m\ge3$ \cite[Theorems 3.1 and 3.2, p.~5]{FI}; our
plethysm family has inner degree three. Ikenmeyer and Panova studied
rectangular Kronecker coefficients, of which $k_n$ is the case of three equal
rectangles; among other things they show that the positivity stretching
factor is at most two for a long first row and that the rectangular Kronecker
semigroup has trivial saturation \cite[Section 5(c),(d), p.~13]{IP}. These are
statements about positivity of the coefficients, not about the indices or the
poles of the stretching functions. The report \cite[Conjecture 3, p.~3]{ICERM}
states the relaxed SH(a) and PH2 for Kronecker coefficients, with the parameter
the largest height, as a conjecture of Mulmuley, and notes that the examples of
\cite{BOR} do not contradict it; it does not state PH3. Our results leave that
conjecture open for Kronecker coefficients.

\paragraph{Kahle and Micha\l{}ek.}
Kahle and Micha\l{}ek~\cite{KM} showed, on the examples $S^3(S^k)$ and
$S^k(S^3)$, that the quasi-polynomials of plethysm need not count lattice points
of inhomogeneous polytopes whose dimension equals their degree. This concerns
the polytope hypothesis PH1 of \cite{GCT6}, not the index hypotheses treated
here.

\paragraph{What remains open.}
(1) The exponent $r$ in Theorem~\ref{thm:kronecker}: whether $K_n$ always has
its pole at the primitive $2^{n-1}$-th roots themselves. (2) The positivity and
saturation indices of the Kronecker family $k_n$. (3) The density clauses
1.6.5(b) and 3.3.2(b). (4) Whether a weakened PH3, with the modular index
replaced by a quantity insensitive to a single pole of exponential order, could
serve the purposes of \cite{GCT6}; the poles found here are simple, isolated in
the two-adic direction, and come from one fixed curve.

\section{Who proved and checked what}\label{sec:provenance}

The Kronecker results (Theorems~\ref{thm:K023}, \ref{thm:K024},
\ref{thm:transfer}, \ref{thm:kronecker}, Lemma~\ref{lem:liezero}) were found
and proved by a run of GPT-6 Astra (OpenAI, ChatGPT Pro mode) on 21--22
November 2026; the support theorem (Section~\ref{sec:support}),
Proposition~\ref{prop:fixedcomponentpole} and the plethysm results
(Sections~\ref{sec:plethysm}--\ref{sec:indices}) and the torus examples by an
earlier run of the same model on 20--21 November. Claude Opus~5.5 at effort max
read every proof of the Kronecker chain, of
Proposition~\ref{prop:fixedcomponentpole} and of Section~\ref{sec:support} and
reran the computations with its own code (22 November); Claude Opus~5.5 at
effort high read the proofs of Theorem~\ref{thm:plethysmPH3}, of the
canonical-algebra identification and of Section~\ref{sec:indices} and recomputed
their numerical claims in small dimensions (21 November). In writing this text
the latter added: the Kempf--Ness proof of closedness for the chain cubic
(replacing a citation), the grading stabilizer of $p_n$ for all $n$ and
Lemma~\ref{lem:orbitclosure} (which give $e_Y=1$ also when $3\mid n$), the
direct count of fixed weights, the route to \eqref{eq:amplitude} through
Proposition~\ref{prop:product}, the Fourier normalization, and the proofs of
Propositions~\ref{prop:torusmodular} and~\ref{prop:torusindices}; it checked
these, and the Hessian locus of Proposition~\ref{prop:chaincubic}(2), by
computation. Claude Fable~5.2 at effort max then read the whole text,
followed every proof, redid the added steps, the constants of
Section~\ref{sec:indices} and Section~\ref{sec:torus} by hand, and tested them
with computations of its own (25 November). No proof in this paper has been
checked by a human.

\begin{thebibliography}{99}

\bibitem{BOR}
E.~Briand, R.~Orellana, M.~Rosas, \emph{Reduced Kronecker coefficients and
counter-examples to Mulmuley's strong saturation conjecture SH}, with an
appendix by K.~Mulmuley, Comput. Complexity \textbf{18} (2009);
arXiv:0810.3163v3 (14 May 2009). Page numbers refer to the arXiv version.

\bibitem{FI}
N.~Fischer, C.~Ikenmeyer, \emph{The computational complexity of plethysm
coefficients}, arXiv:2002.00788v1 (3 February 2020).

\bibitem{ICERM}
J.~A.~Grochow, K.~Rusek, \emph{Report on ``Mathematical aspects of P vs.\ NP
and its variants''}, workshop at ICERM, 1--5 August 2011 (organizers S.~Basu,
J.~M.~Landsberg, J.~M.~Rojas), arXiv:1203.2888v2 (21 March 2012).

\bibitem{IMW}
C.~Ikenmeyer, K.~D.~Mulmuley, M.~Walter, \emph{On vanishing of Kronecker
coefficients}, arXiv:1507.02955v2 (21 July 2017).

\bibitem{IP}
C.~Ikenmeyer, G.~Panova, \emph{Rectangular Kronecker coefficients and
plethysms in geometric complexity theory}, arXiv:1512.03798v2 (11 August 2017).

\bibitem{KM}
T.~Kahle, M.~Micha\l{}ek, \emph{Obstructions to combinatorial formulas for
plethysm}, Electron. J. Combin. \textbf{25}(1) (2018), \#P1.41;
arXiv:1507.07131v3.

\bibitem{Losev}
I.~Losev, \emph{The Kempf--Ness theorem and invariant theory},
arXiv:math/0605756.

\bibitem{Luna}
D.~Luna, \emph{Slices \'etales}, Bull. Soc. Math. France, M\'emoire \textbf{33}
(1973), 81--105.

\bibitem{Macdonald}
I.~G.~Macdonald, \emph{Symmetric functions and Hall polynomials}, 2nd ed.,
Oxford University Press, 1995.

\bibitem{GCT6}
K.~D.~Mulmuley, \emph{Geometric complexity theory VI: the flip via saturated
and positive integer programming in representation theory and algebraic
geometry}, arXiv:0704.0229v4 (22 January 2009). Page numbers refer to this
version.

\bibitem{Stacks}
The Stacks Project Authors, \emph{The Stacks Project},
\url{https://stacks.math.columbia.edu}, Section 15.29 (Tag 0621) and
Lemma 15.31.2 (Tag 062F).

\end{thebibliography}

\end{document}
